\documentclass[11pt]{article}
\usepackage[T1]{fontenc}
\usepackage{microtype}
\usepackage{amssymb,amsthm,mathtools}
\usepackage[margin=1in]{geometry}
\usepackage[unicode=true,colorlinks=true,citecolor=blue,linkcolor=blue,urlcolor=blue]{hyperref}
\usepackage{bookmark}
\hypersetup{
  pdftitle={A Floor-Saving Quadratic-over-Logarithmic Lower Bound for Representation Tops in Erd\H{o}s Problem \#1194},
  pdfauthor={GPT-5.5 Pro and Lech Mazur},
  pdfsubject={Perfect difference sets and representation tops},
  pdfkeywords={Erd\H{o}s Problem 1194, perfect difference sets, Sidon sets, representation functions}
}
\allowdisplaybreaks[1]

\newtheorem{theorem}{Theorem}
\newtheorem{lemma}{Lemma}
\newcommand{\N}{\mathbb{N}}
\newcounter{proofheading}
\newcounter{proofsubheading}[proofheading]
\newcommand{\proofheading}[1]{%
  \refstepcounter{proofheading}%
  \medskip\phantomsection%
  \pdfbookmark[1]{#1}{proofheading.\theproofheading}%
  \noindent\textbf{#1}\par\smallskip}
\newcommand{\proofsubheading}[1]{%
  \refstepcounter{proofsubheading}%
  \smallskip\phantomsection%
  \pdfbookmark[2]{#1}{proofsubheading.\theproofheading.\theproofsubheading}%
  \noindent\textit{#1}\par\smallskip}

\title{A Floor-Saving Quadratic-over-Logarithmic Lower Bound for Representation Tops in Erd\H{o}s Problem \#1194}
\author{GPT-5.5 Pro\thanks{Large language model by OpenAI; performed nearly all of the mathematical derivation and LaTeX drafting through prompted interaction, with human prompting, review, and editing.}
\and
Lech Mazur\thanks{Human contributor, prompter, curator, editor, auditor, and responsible contact for public posting.}}
\date{}

\begin{document}
\maketitle

\begin{abstract}
Erd\H{o}s Problem \#1194 asks how fast the representation top must grow when
$A\subseteq\mathbb{N}=\{1,2,3,\ldots\}$ is a perfect difference set in the infinite-integer sense used here,
meaning that every integer $n\ge1$ has a unique ordered representation
\[
n=u(n)-v(n),
\qquad
u(n),v(n)\in A,
\qquad
u(n)>v(n).
\]
Thus $u(n)$ is the larger element in the representation of $n$, corresponding to the top entry
in the original problem statement, not the $n$th element of $A$. The main result proves an
explicit quadratic-over-logarithmic lower bound in this setting. Let
\[
\lambda=\frac1{2\log 2},
\qquad
B_1=\frac12+\gamma-\log(\log 2),
\]
where $\gamma$ is the Euler--Mascheroni constant. All logarithms in the displayed formulae are
natural. For every fixed $B>B_1$, there are arbitrarily large $n$ for which the denominator below is
positive and
\[
u(n)>\lambda\,\frac{n^2}{\log n-\log\log n+B}.
\]
Equivalently, for this fixed $B$, the eventual bound
$u(n)\le\lambda n^2/(\log n-\log\log n+B)$ cannot hold. Consequently, for every fixed real $C$,
\[
 \limsup_{n\to\infty}
 \frac{u(n)(\log n-\log\log n+C)}{n^2}
 \ge \frac1{2\log 2}.
\]
A companion Lean proof is available at
\url{https://github.com/lechmazur/erdos_1194/}. No pointwise infinitely-often inequality at
$B=B_1$ is asserted.

\end{abstract}

\medskip
\noindent\textbf{MSC 2020.} 11B13; 11B34; 05B10.

\noindent\textbf{Keywords.} perfect difference sets; Sidon sets; Erd\H{o}s problems; representation functions.

All logarithms are natural. Let
\[
\N=\{1,2,3,\ldots\}.
\]
Let
\[
\lambda:=\frac1{2\log 2}.
\]
Let $\gamma$ denote the Euler--Mascheroni constant,
\[
\gamma=\lim_{K\to\infty}\left(\sum_{m=1}^K\frac1m-\log K\right),
\]
and put
\[
B_1:=\frac12+\gamma-\log(\log 2).
\]
Numerically,
\[
\lambda=0.721347520444\ldots,
\qquad
B_1=1.443728585483\ldots.
\]
Since $0<\log 2<1$, we have $\log(\log 2)<0$; in particular, $B_1>0$.

\medskip
\noindent\textit{Asymptotic notation.} Throughout, $F\ll_{\mathcal P}G$ means
$|F|\le C_{\mathcal P}G$ on the stated range, $F\gg_{\mathcal P}G$ means
$G\ll_{\mathcal P}F$, and $F\asymp_{\mathcal P}G$ means both estimates hold. The same
subscript convention is used for $O$-terms.

\medskip
\noindent\textit{Context and status.} The hypothesis that every $n\ge1$ has a unique ordered
representation $n=a-b$ with $a,b\in A$ and $a>b$ appears in Erd\H{o}s's 1980 problem survey
\cite[p.~100]{Erdos1980}; it is also catalogued as Erd\H{o}s Problem \#1194
\cite{Bloom1194}. Here, sets satisfying this hypothesis are called perfect difference sets in
this infinite-integer setting; this terminology is local to the present problem and is not intended
to invoke the finite cyclic perfect-difference-set convention. Recent work on the finite
Sidon-extension conjecture, including counterexamples to that finite conjecture, is logically
independent of the infinite representation-top growth function studied here
\cite{AlexeevMixon2026}. Throughout, write the unique representation of $n$ as
\[
n=u(n)-v(n),
\qquad
u(n),v(n)\in A,
\qquad
u(n)>v(n).
\]
Then $u(n)$ is the representation top, corresponding to the larger element in the original
problem statement, not the $n$th element of $A$.

The catalogue page records a greedy construction due to Lev, with growth $u(n)\ll n^3$ in this
representation-top notation \cite{Lev2004,Bloom1194}. It also records the standard Sidon-density
argument, based on Erd\H{o}s's upper bound for infinite Sidon sequences as presented, for example,
by Halberstam and Roth, which yields $u(n)\gg n\log n$ infinitely often
\cite{HalberstamRoth1966,Bloom1194}. The discussion page also records subsequent progress around
this problem: Liam Price linked a GPT-5.4 Pro note, and Thomas Bloom distilled from the ensuing
discussion an argument giving $a_n\gg n^{3/2}$ infinitely often. Bloom then reported a GPT Pro
improvement to $a_n\gg n^{2-o(1)}$ infinitely often, and suggested an optimized
$n^2/f(n)$-type form under the convergence condition $\sum_n 1/(n f(n))<\infty$
\cite{Bloom1194}. The present draft gives a separate fully written logarithmic lower bound with
explicit constants and a companion Lean formalization.

The perfect-difference hypothesis implies the Sidon property that all unordered pair sums
$a+b$, with $a\le b$ and $a,b\in A$, are distinct. Indeed, suppose $a\le b$, $c\le d$, and
$a+b=c+d$. If $b=d$, then $a=c$. Otherwise, after swapping the two pairs if necessary, $b>d$, and
then
\[
b-d=c-a>0.
\]
The two ordered pairs cannot coincide: if $(b,d)=(c,a)$, then $b=c$ and $d=a$, so
$b>d$ gives $c>a$, while $c\le d=a$ gives $c\le a$. This gives two distinct ordered
representations, $(b,d)$ and $(c,a)$, of the same positive difference, contradicting uniqueness.
If all representation tops up to difference $n$ were too
small, the first $n$ unique difference pairs would force too many elements of $A$ inside a bounded
interval, contradicting the classical infinite-Sidon counting bound at that scale. Cilleruelo and
Nathanson constructed dense perfect difference sets from dense Sidon sets, showing that the
existence theory for perfect difference sets is linked to the density theory of Sidon sets
\cite{CillerueloNathanson2008}.

\medskip
\noindent\textit{Scope.} Existence is known from the cited constructions. The theorem is universal
over all sets satisfying the hypothesis; it does not use or reprove any construction of such sets.
The proof starts with an arbitrary set $A$ satisfying the perfect-difference hypothesis and derives
the stated lower bound for its representation tops. This draft does not determine the optimal
growth of $u(n)$.

\begin{theorem}\label{thm:main}
Let $A\subseteq\N$ be such that, for every $n\ge 1$, there is exactly one ordered pair
$(a,b)\in A^2$ with $a>b$ and $a-b=n$. Write this pair as $(u(n),v(n))$, so
$n=u(n)-v(n)$ and $u(n)>v(n)$.
For every fixed $B>B_1$, there are arbitrarily large integers $n$ such that
\[
 u(n)>
 \lambda\,
 \frac{n^2}{\log n-\log\log n+B}.
\]
No uniformity as $B\downarrow B_1$ is asserted. The displayed denominator is positive for all
sufficiently large $n$, so there is no large-$n$ sign ambiguity.

Equivalently, for this fixed $B$, the eventual non-strict bound
$u(n)\le \lambda n^2/(\log n-\log\log n+B)$ cannot hold.

As a limsup consequence, for every fixed real $C$,
\[
 \limsup_{n\to\infty}
 \frac{u(n)(\log n-\log\log n+C)}{n^2}
 \ge
 \frac1{2\log 2}.
\]
Taking $C=B_1$ gives the endpoint limsup version. Since
$\log n/(\log n-\log\log n)\to1$, the conclusion with $C=0$ also gives
\[
 \limsup_{n\to\infty}\frac{u(n)\log n}{n^2}
 \ge \lambda.
\]
\end{theorem}

\noindent\textit{Remark.} Theorem~\ref{thm:main} does not determine the optimal growth of $u(n)$.
It is a universal lower bound for representation tops, not a construction and not a matching
upper bound. The proof below does not establish the pointwise infinitely-often inequality at the
endpoint $B=B_1$. The endpoint statement is obtained only by applying the proved result with any fixed
$B^*>B_1$ and then taking a limsup along the resulting unbounded sequence; the proof never
substitutes $B=B_1$ into the contradiction argument. More generally, any fixed additive constant
in the limsup denominator gives the same leading constant. No optimality of $B_1$ as a sharp
pointwise threshold is claimed.

\noindent\textit{Proof roadmap.} The proof argues by contradiction. The assumed upper bound
gives a spacing rule for elements of $A$, which yields a continuous majorant for the counting
identity by tops. The integer floor in that identity produces a fixed normalized saving, which
contradicts the continuous asymptotic when $B>B_1$. More precisely, the continuous second-order
term is reduced by the floor-saving term $2\lambda(1-\gamma)X/\log X$, where
$\{q\}$ denotes the fractional part of $q$; the integral
$1-\gamma=\int_1^\infty \{q\}q^{-2}\,dq>0$ appears after the later change of variables
$q=X/\psi(t)$. In compressed form, the dependency chain is
\[
\text{spacing}\Rightarrow\text{counting majorant}\Rightarrow J_B(X)-E(X)
\Rightarrow\text{a negative }X/\log X\text{ coefficient}.
\]

\noindent\textit{Constant bookkeeping.} The final comparison uses
\[
\lambda=\frac1{2\log 2},
\qquad
B_1=\frac12+\gamma-\log(\log 2),
\]
and the two second-order coefficients
\[
C_{\rm cont}(B):=\lambda\bigl(3-2B-2\log(\log 2)\bigr),
\qquad
C_{\rm floor}:=2\lambda(1-\gamma).
\]
At the endpoint,
\[
C_{\rm cont}(B_1)=C_{\rm floor}=0.609948863612\ldots.
\]
The continuous coefficient minus the floor-saving coefficient is
\[
C_{\rm cont}(B)-C_{\rm floor}
=\lambda\bigl(3-2B-2\log(\log 2)\bigr)-2\lambda(1-\gamma)
=-2\lambda(B-B_1).
\]
The main auxiliary objects below are the eventual majorant $\Phi$, its inverse $\psi$, the
nonincreasing function $f=1/\psi$ on the inverse tail, the top-window integral $G_X$, the
continuous majorant $J_B(X)$, and the floor saving $E(X)$.

\begin{proof}
Use the notation $u(n),v(n)$ from the theorem. Fix
\[
B>B_1.
\]
All thresholds and implicit constants below are attached to the fixed set $A$, the fixed parameter
$B$, and the contradiction data introduced below; no uniformity over all admissible sets $A$, or
as $B\downarrow B_1$, is asserted.
Recall that
\[
\lambda=\frac1{2\log 2}.
\]
Since $B>B_1$ and $B_1>0$, we have $B>0$. Put
\[
H(w):=w-\log w+B\qquad(w>0).
\]
For $w>0$, the elementary inequality $w-\log w\ge1$, together with $B>B_1>0$, gives
\[
H(w)>0.
\]
Also
\[
H'(w)=1-\frac1w,
\]
and
\[
2H(w)-H'(w)
=
2(w-\log w)+2B-1+\frac1w
\ge
1+2B+\frac1w>0.
\]
Thus, for every $x>1$, define
\[
\Phi(x):=\lambda\frac{x^2}{\log x-\log\log x+B}
=\lambda\frac{x^2}{H(\log x)}.
\]
Then
\[
\Phi(e^w)=\lambda\frac{e^{2w}}{H(w)}
\qquad(w>0),
\]
and
\[
\frac{d}{dw}\Phi(e^w)
=
\lambda e^{2w}\frac{2H(w)-H'(w)}{H(w)^2}>0.
\]
Also, as $x\to\infty$,
\[
\Phi(x)\sim \lambda\frac{x^2}{\log x}\to\infty.
\]
Assume, for contradiction, that
\[
u(n)\le \Phi(n)
\tag{1}\label{eq:upper-bound-assumption}
\]
for all sufficiently large $n$.

Choose and fix an integer $N_0\ge 3$ large enough that
\[
u(n)\le\Phi(n)\qquad(n\ge N_0),
\]
and that
\[
\frac12\log x\le H(\log x)\le 2\log x,
\qquad
\log x\le 2H(\log x)-H'(\log x)\le 4\log x
\qquad(x\ge N_0).
\]
Such an $N_0$ exists because the assumed bound holds eventually and
\[
H(\log x)\sim\log x,
\qquad
2H(\log x)-H'(\log x)\sim2\log x
\qquad(x\to\infty).
\]
The displayed comparisons are used only for tail estimates; positivity of $H(\log x)$ and
$2H(\log x)-H'(\log x)$ already holds for every $x>1$.
Since
\[
\Phi'(x)
=
\lambda x\frac{2H(\log x)-H'(\log x)}{H(\log x)^2},
\]
we have $\Phi'(x)>0$ for every $x\ge N_0$. Let
\[
T_*:=\Phi(N_0).
\]
Since $\lambda>0$, $N_0>0$, and $H(\log N_0)>0$, the definition
$T_*=\lambda N_0^2/H(\log N_0)$ gives $T_*>0$. Since $\Phi$ is continuous, strictly increasing on
$[N_0,\infty)$, and tends to infinity, it maps $[N_0,\infty)$ bijectively onto $[T_*,\infty)$.
Let $\psi$ be the inverse of $\Phi$ on this interval. Since $\Phi$ is $C^1$ and $\Phi'>0$ on
$[N_0,\infty)$, the inverse $\psi$ is $C^1$ on $(T_*,\infty)$. All later uses of $\psi$ occur
at arguments at least $T_*$, after the relevant lower thresholds have been imposed. Thus
whenever $\psi$ appears below, it denotes this branch only; inversions of $\Phi$ are never
taken outside $[N_0,\infty)$.

Define
\[
f(t):=\frac1{\psi(t)}\qquad(t\ge T_*).
\]
Extend $f$ to $[0,T_*)$ by
\[
f(t):=f(T_*)\qquad(0\le t<T_*).
\]
Then $f$ is positive, nonincreasing, locally bounded, and measurable on $[0,\infty)$. Indeed,
$\psi$ is increasing on $[T_*,\infty)$, so $1/\psi$ is nonincreasing on $[T_*,\infty)$. Also
$\psi(T_*)=N_0$, and hence $f(t)\le f(T_*)$ for $t\ge T_*$. Therefore the constant extension to
$[0,T_*)$ preserves monotonicity across the joining point. Moreover, $f$ is $C^1$ on the smooth
tail $(T_*,\infty)$. All integrals below are Lebesgue integrals.

The artificial extension on $[0,T_*)$ is only a bookkeeping device. Only the values $t\ge T_*$
are tied to the inverse branch $\psi$; values below $T_*$ are used solely inside integrals. Since
$T_*$ is fixed and $f$ is bounded there, every contribution from this interval is $O_{A,B}(1)$,
hence negligible on the scales $\sqrt{X/\log X}$ and $X/\log X$ used later. It is also negligible
in the normalized $L^1$ rescaling used below, where its contribution is $O_{A,B}(\sqrt{Y/X})$.

\begin{lemma}[Uniform inverse-branch convention on compact $q$-ranges]
For every fixed $Q\ge1$, all sufficiently large $X$ satisfy $X/Q\ge N_0$, and then
\[
\psi(\Phi(X/q))=X/q
\qquad(1\le q\le Q).
\]
Consequently, after increasing the lower threshold in $X$ when necessary, every later occurrence
of $\psi$ is evaluated on the inverse branch from $[T_*,\infty)$ to $[N_0,\infty)$ fixed above.
\end{lemma}

\begin{proof}
For fixed $Q\ge1$, the inequality $X/Q\ge N_0$ holds for all sufficiently large $X$. Then
$X/q\ge X/Q\ge N_0$ for every $1\le q\le Q$, and the identity follows from the definition of the
inverse branch. The displayed choices used later, including $M=\Phi(X)$, $\Phi(X/Q)$, and
$\Phi(X/q)$ for fixed $Q$ and $1\le q\le Q$, are covered by this compact-range observation;
separately, $X\ge T_*$ for all sufficiently large $X$.
\end{proof}

\proofheading{1. Spacing}

Write
\[
A=\{x_1<x_2<x_3<\cdots\}.
\]
This enumeration is infinite because every positive integer occurs as a difference of two
elements of $A$. Hence $A$ is unbounded.

Since $N_0\ge3$, the finite maximum below is over a nonempty set. Choose an integer
\[
Z_0\in\mathbb N,
\qquad
Z_0>\max\left(T_*,\max_{1\le d<N_0}u(d)\right).
\]
The integrality of $Z_0$ is used only when bounding integer tops in the short interval
$(X,X+Z_0]$.

\medskip
\noindent\textit{Threshold and uniformity convention.} Lower thresholds may be enlarged finitely
many times after the fixed data $A,B,N_0,T_*,Z_0$ and any fixed truncation parameter have been
chosen. Thus all later occurrences of $\log\log$, $\Phi$, $\psi$, and denominators of the form
$\log z-\log\log z+C$ are on ranges where they are defined and have the asserted sign.

Unless explicitly stated otherwise, all $O$- and $o$-terms are taken after the contradiction data
have been fixed. Constants depending on $A,B,N_0,T_*,Z_0$ are written as $O_{A,B}$; lower
thresholds may also depend on previously fixed auxiliary parameters. If a truncation parameter such
as $U$ or $Q$ appears, it is fixed before $X\to\infty$ and is sent to its limiting value only after
that limit. Any tail estimate later used after $U$ or $Q$ is sent to infinity states the needed
independence of the displayed implicit constant; no truncation parameter varies with $X$. Whenever
$Y=\log X$ is used below, the scale separations $X/Y\to\infty$ and $X/Y^2\to\infty$ are also used
only on this eventual range.

The notation convention stated above remains in force. Extra subscripts record extra fixed data
on which the implicit constant may depend.

\proofsubheading{Spacing and gap-integral domination}

\begin{lemma}[Spacing and gap-integral domination]
If $y<z$ are elements of $A$ and $z>Z_0$, then
\[
z-y\ge \psi(z).
\tag{2}\label{eq:spacing}
\]
Consequently,
\[
\int_y^z f(s)\,ds\ge1.
\tag{3}\label{eq:gap-integral}
\]
\end{lemma}

\begin{proof}
Put
\[
d:=z-y.
\]
Then $d$ is a positive integer.
The ordered pair $(z,y)$ represents $d$. By uniqueness of the representing pair for $d$,
\[
u(d)=z.
\]
If $d<N_0$, then
\[
z=u(d)\le \max_{1\le m<N_0}u(m)<Z_0,
\]
contradicting $z>Z_0$. Hence $d\ge N_0$. By the assumed upper bound
\eqref{eq:upper-bound-assumption},
\[
z=u(d)\le \Phi(d).
\]
Since $d\ge N_0$, $z>Z_0>T_*$, and $\Phi$ is increasing on $[N_0,\infty)$, applying the
inverse branch $\psi:[T_*,\infty)\to[N_0,\infty)$ to $z\le\Phi(d)$ gives
\[
d\ge \psi(z).
\]
Thus $z-y=d\ge\psi(z)$, proving \eqref{eq:spacing}.

Since $f$ is nonincreasing, $f(s)\ge f(z)$ for $y\le s\le z$. Also $z>Z_0>T_*$ gives
$f(z)=1/\psi(z)$. Hence
\[
\int_y^z f(s)\,ds
\ge
(z-y)f(z)
=
\frac{z-y}{\psi(z)}
\ge1.
\]
\end{proof}

\proofsubheading{Counting-function majorant}

For real $T\ge0$, define
\[
\mathcal N_A(T):=|A\cap[1,T]|,
\]
where $[1,T]$ is understood to be empty if $0\le T<1$.

\begin{lemma}[Counting-function majorant]
For all real $T\ge0$,
\[
\mathcal N_A(T)\le C_*+\int_0^T f(s)\,ds
\tag{4}\label{eq:counting-function-bound}
\]
for some constant $C_*=C(A,B,N_0,T_*,Z_0)$.
\end{lemma}

\begin{proof}
In the argument below one may take $C_*=j_0$ after $j_0$ has been defined. By unboundedness and
the well-ordering of $\N$, there is a first element of $A$ exceeding $Z_0$; call it $x_{j_0}$. If
$\mathcal N_A(T)\le j_0$, then
\eqref{eq:counting-function-bound} holds with $C_*\ge j_0$. Otherwise, let
$m=\mathcal N_A(T)>j_0$.
Then $x_m\le T$, and by \eqref{eq:gap-integral},
\[
m-j_0
\le
\sum_{j=j_0+1}^m\int_{x_{j-1}}^{x_j}f(s)\,ds
\le
\int_0^T f(s)\,ds.
\]
Thus
\[
\mathcal N_A(T)=m\le j_0+\int_0^T f(s)\,ds,
\]
which proves \eqref{eq:counting-function-bound} with $C_*=j_0$.
\end{proof}

\proofheading{2. Estimates for \texorpdfstring{$f$}{f}}

For the tail estimates in this part, let $t>T_*$ and write
\[
r=\psi(t).
\]
Then
\[
t=\Phi(r)=\lambda\frac{r^2}{\log r-\log\log r+B}=\lambda\frac{r^2}{H(\log r)}.
\]
As $t\to\infty$, also $r\to\infty$. Since
\[
H(\log r)\asymp_B \log r
\]
for large $r$, we have
\[
\log t
=
2\log r-\log H(\log r)+O_B(1)
=
2\log r+O_B(\log\log r).
\]
Since $O_B(\log\log r)=o(\log r)$, this first implies
\[
\log r\asymp_B \log t.
\]
Consequently $\log H(\log r)=O_B(\log\log t)$, because $H(\log r)\asymp_B\log r$ on the tail.
Substituting this back into
\[
\log t=2\log r-\log H(\log r)+O_B(1)
\]
gives
\[
\log r=\frac12\log t+O_B(\log\log t).
\]
Equivalently,
\[
\log r=\frac12\log t
\left(1+O_B\left(\frac{\log\log t}{\log t}\right)\right).
\]
Therefore
\[
\log\log r
=
\log\log t+
\log\left(\frac12+O_B\left(\frac{\log\log t}{\log t}\right)\right)
=\log\log t+O_B(1).
\]
Also,
\[
f(t)^2
=
\frac1{r^2}
=
\frac{\lambda}{t\,H(\log r)}.
\]
Because
\[
H(\log r)
=\log r-\log\log r+B
=\frac12\log t+O_B(\log\log t),
\]
we have
\[
\frac1{H(\log r)}
=
\frac2{\log t}
\left(1+O_B\left(\frac{\log\log t}{\log t}\right)\right).
\]
Together with $f(t)^2=\lambda/(tH(\log r))$, and using that the factor in parentheses is
$1+o(1)>0$ on the tail and $\sqrt{1+O(\varepsilon)}=1+O(\varepsilon)$ for small
$\varepsilon$, this gives after taking the positive square root, uniformly for all sufficiently
large $t$,
\[
f(t)
=
\sqrt{\frac{2\lambda}{t\log t}}
\left(1+O_B\left(\frac{\log\log t}{\log t}\right)\right).
\tag{5}\label{eq:f-asymptotic}
\]
In particular,
\[
f(t)\ll_B \frac1{\sqrt{t\log t}}
\qquad(t\to\infty).
\tag{6}\label{eq:f-upper-bound}
\]

\proofsubheading{Integral asymptotic for \texorpdfstring{$f$}{f}}

\begin{lemma}[Integral asymptotic for \texorpdfstring{$f$}{f}]
As $X\to\infty$,
\[
\int_0^X f(t)\,dt
=
2\sqrt{\frac{2\lambda X}{\log X}}
+
o\left(\sqrt{\frac X{\log X}}\right).
\tag{7}\label{eq:f-integral-asymptotic}
\]
\end{lemma}

\begin{proof}
We first record the elementary estimate
\[
\int_{T_1}^V\frac{dt}{\sqrt{t\log t}}
\ll
\sqrt{\frac V{\log V}}
\qquad(V\to\infty),
\tag{8}\label{eq:elementary-integral-bound}
\]
for any fixed sufficiently large $T_1$. For $V\ge e^2$,
\[
\frac{d}{dV}\sqrt{\frac V{\log V}}
=
\frac{\log V-1}{2\sqrt V(\log V)^{3/2}}
\asymp
\frac1{\sqrt{V\log V}},
\]
so increasing $T_1$ if necessary and integrating gives \eqref{eq:elementary-integral-bound}.

Now put
\[
Y:=\log X.
\]
Since $X/Y^2\to\infty$, all thresholds below are exceeded for sufficiently large $X$. The
bounded initial interval contributes $O_{A,B}(1)$. By \eqref{eq:f-upper-bound}, for any fixed
sufficiently large $T_1$,
\[
\int_0^{X/Y^2} f(t)\,dt
=
O_{A,B}(1)+O_B\left(\int_{T_1}^{X/Y^2}\frac{dt}{\sqrt{t\log t}}\right).
\]
Using \eqref{eq:elementary-integral-bound}, the last integral is
\[
\ll_B
\sqrt{\frac{X/Y^2}{\log(X/Y^2)}}
\ll_B
\frac{\sqrt X}{Y^{3/2}}
=
o\left(\sqrt{\frac X Y}\right).
\]
Thus
\[
\int_0^{X/Y^2} f(t)\,dt=o\left(\sqrt{\frac X Y}\right).
\]
On the interval $X/Y^2\le t\le X$, we have
\[
\log t=Y+O(\log Y)
\]
uniformly. Hence the error term in \eqref{eq:f-asymptotic} is uniformly $o(1)$ on this
interval; explicitly,
\[
\sup_{X/Y^2\le t\le X}
\left|
\frac{f(t)}{\sqrt{2\lambda/(tY)}}-1
\right|\to0.
\]
Therefore
\[
f(t)=\sqrt{\frac{2\lambda}{Y}}\,t^{-1/2}(1+o(1))
\]
uniformly for $X/Y^2\le t\le X$. Also,
\[
\int_{X/Y^2}^X t^{-1/2}\,dt
=
2\sqrt X\left(1-\frac1Y\right)
=
2\sqrt X+O\left(\frac{\sqrt X}{Y}\right).
\]
It follows that
\[
\int_{X/Y^2}^X f(t)\,dt
=
2\sqrt{\frac{2\lambda X}{Y}}
+
o\left(\sqrt{\frac X Y}\right),
\]
which proves \eqref{eq:f-integral-asymptotic}.
\end{proof}

Since $Z_0$ is fixed, the same estimate remains
valid with $X+Z_0$ in place of $X$. Indeed, apply \eqref{eq:f-integral-asymptotic} with
$X+Z_0$ in place of $X$; writing $Y=\log X$,
\[
\sqrt{\frac{X+Z_0}{\log(X+Z_0)}}
\sim
\sqrt{\frac X{\log X}}
=
\sqrt{\frac X Y}.
\]
Thus
\[
\int_0^{X+Z_0} f(t)\,dt
=
2\sqrt{\frac{2\lambda X}{Y}}
+
o\left(\sqrt{\frac X Y}\right)
=
O_{A,B}\left(\sqrt{\frac X Y}\right).
\]
Together with \eqref{eq:counting-function-bound}, this gives
\[
\mathcal N_A(X+Z_0)=O_{A,B}\left(\sqrt{\frac X Y}\right).
\tag{9}\label{eq:N-XplusZ0-bound}
\]

Finally, on the smooth tail $t>T_*$,
\[
|f'(t)|\ll_B \frac{f(t)}t.
\tag{10}\label{eq:f-derivative-bound}
\]
Indeed, if $t=\Phi(r)$, then $r=\psi(t)\ge N_0$ and $f(t)=1/r$, so
\[
f'(t)=-\frac1{r^2\Phi'(r)}.
\]
Now
\[
\Phi'(r)=\lambda r\frac{2H(\log r)-H'(\log r)}{H(\log r)^2}.
\]
Therefore
\[
|f'(t)|
=
\frac{H(\log r)^2}
{\lambda r^3\bigl(2H(\log r)-H'(\log r)\bigr)}.
\]
Since
\[
H(w)\sim w,
\qquad
2H(w)-H'(w)\sim2w,
\]
and since the displayed bounds imposed when choosing $N_0$ apply with $x=r$ for every $r\ge
N_0$, these comparisons hold with uniform upper and lower constants. Hence
\[
\Phi'(r)\asymp_B \frac r{H(\log r)},
\]
and the exact identity above gives
\[
|f'(t)|\asymp_B\frac{H(\log r)}{r^3}.
\]
On the other hand,
\[
t\asymp_B\frac{r^2}{H(\log r)},
\qquad
f(t)=\frac1r.
\]
Therefore
\[
|f'(t)|\ll_B\frac{f(t)}t.
\]
\begin{lemma}[Local comparability and Lipschitz control]
For every fixed $c\in(0,1)$, there is a threshold $T(c)>T_*/c$ such that, whenever
$t\ge T(c)$ and $ct\le s\le t$,
\[
f(s)\asymp_{B,c} f(t),
\tag{11}\label{eq:local-comparability}
\]
and
\[
|f(s)-f(t)|\ll_{B,c}\frac{f(t)}t\,|t-s|.
\tag{12}\label{eq:local-lipschitz}
\]
For fixed $c$, the implicit constants here depend only on $B$ and $c$, not on the later
truncation parameters $U$ or $Q$.
\end{lemma}

\begin{proof}
If $ct\le s\le t$, then
\[
\log s=\log t+O_c(1),
\]
so $s^{-1/2}\asymp_c t^{-1/2}$ and $\log s\asymp_c\log t$. Applying
\eqref{eq:f-asymptotic} to both $s$ and $t$, after increasing $T(c)$ if necessary, gives
\eqref{eq:local-comparability}. For $t\ge T(c)$, the full interval $[s,t]$ lies in the smooth
tail $(T_*,\infty)$. If $\xi\in[s,t]$, then $ct\le \xi\le t$, so
\eqref{eq:local-comparability}, applied with $\xi$ in place of $s$, gives
\[
f(\xi)\asymp_{B,c}f(t).
\]
By \eqref{eq:f-derivative-bound},
\[
|f'(\xi)|\ll_B \frac{f(\xi)}{\xi}
\ll_{B,c}\frac{f(t)}t.
\]
The mean-value theorem now gives \eqref{eq:local-lipschitz}.
\end{proof}

\proofheading{3. Counting by tops}

This section records the exact top-counting identity used in the contradiction.

From now on, $X$ tends to infinity through positive integers in counting identities. The same
symbol $X$ is also used as a real analytic parameter, for example as an endpoint of integration
or as an argument of $\Phi$ and $\psi$.

Let $X\to\infty$ through positive integers and put
\[
Y:=\log X,
\qquad
M:=\Phi(X).
\]
The quantity $M$ need not be an integer. Sums over $t\in A$ with $t\le M$ are sums over integer
elements of $A$ bounded by the real number $M$. The analytic estimates are applied to the same
integer values of $X$, with $X$ also used as a real parameter inside integrals. Throughout the
remaining proof, $X$ is assumed large enough for every analytic interval whose nonemptiness is
required to be nonempty and for all previously fixed thresholds to be exceeded. Discrete sets such
as $A\cap(X+Z_0,M]$ may of course be empty and are handled separately.

For all sufficiently large integer $X$, every $1\le n\le X$ has
\[
u(n)\le M.
\]
Indeed, for the finitely many $n<N_0$, each $u(n)$ is fixed and $M=\Phi(X)\to\infty$, so this
holds once $X$ is large enough. For $n\ge N_0$, the assumed bound
\eqref{eq:upper-bound-assumption} gives
\[
u(n)\le\Phi(n)\le\Phi(X)=M,
\]
because $n\le X$ and $\Phi$ is increasing on $[N_0,\infty)$.

Also,
\[
M=\Phi(X)\asymp_B\frac{X^2}{Y},
\]
so $M>X+Z_0$ for all sufficiently large $X$, since $Z_0$ is fixed. Since every difference
$1\le n\le X$ has top at most $M$, counting differences by their top gives the following exact
identity.

\begin{lemma}[Top-counting identity]
For all sufficiently large integer $X$, the exact top-counting identity is
\[
X=
\sum_{\substack{t\in A\\ t\le M}}
\#\{b\in A:t-X\le b<t\}.
\tag{13}\label{eq:top-counting-identity}
\]
\end{lemma}

\begin{proof}
Each $1\le n\le X$ contributes its unique representing pair
\[
n=u(n)-v(n),
\]
whose top satisfies $u(n)\le M$, so the map from counted pairs to their differences is
surjective onto $\{1,2,\ldots,X\}$. Conversely, every counted pair $(t,b)$ gives a positive
difference
\[
t-b\in[1,X].
\]
Uniqueness ensures that two counted pairs cannot give the same difference, so this map is
injective on the counted pairs. Therefore the count is exactly $X$.
\end{proof}

No loss has been introduced in \eqref{eq:top-counting-identity}; all later inequalities enter
only after splitting the exact count. We split the right-hand side according as $t\le X$ or
$t>X$.

\proofheading{4. Small tops}

Since $M>X$, the portion of \eqref{eq:top-counting-identity} with $t\le X$ is exactly
\[
\binom{\mathcal N_A(X)}2.
\]
Each two-element subset of $A\cap[1,X]$ contributes exactly one ordered pair, with the larger
element as top and the smaller element as bottom. Indeed, if $t\le X$, then every $b\in A$ with
$b<t$ gives a difference
\[
1\le t-b\le X.
\]
Conversely, since $A\subseteq\N$, every such $b$ satisfies $b\ge1$; and since $t-X\le0$, the
condition $t-X\le b<t$ in \eqref{eq:top-counting-identity} is automatic for every $b\in A$ with
$b<t$. By \eqref{eq:counting-function-bound}, there is a fixed constant
$C_{\rm sp}=C_{\rm sp}(A,B,N_0,T_*,Z_0)$ such that
\[
\mathcal N_A(X)\le F(X)+C_{\rm sp},
\qquad
F(X):=\int_0^X f(t)\,dt.
\]
The artificial extension of $f$ on $[0,T_*)$ changes $F(X)$ only by a fixed $O_{A,B}(1)$
amount relative to the smooth tail. The discarded term in
\[
\binom{\mathcal N_A(X)}2=\frac12\mathcal N_A(X)^2-\frac12\mathcal N_A(X)
\]
is harmless at the present scale, because $\mathcal N_A(X)\le \mathcal N_A(X+Z_0)=O_{A,B}(\sqrt{X/Y})$ by
\eqref{eq:N-XplusZ0-bound}, and
\[
\sqrt{X/Y}=o(X/Y).
\]
Therefore the cruder bound
\[
\binom{\mathcal N_A(X)}2
\le
\frac12\mathcal N_A(X)^2
\le
\frac12(F(X)+C_{\rm sp})^2
\]
is sharp enough for the second-order comparison.
Expanding,
\[
\frac12(F(X)+C_{\rm sp})^2
=
\frac12 F(X)^2+O_{A,B}(F(X))+O_{A,B}(1).
\]
By \eqref{eq:f-integral-asymptotic},
\[
F(X)=2\sqrt{\frac{2\lambda X}{\log X}}+o\left(\sqrt{\frac X{\log X}}\right),
\]
and no second-order expansion of $F$ is needed. Hence
\[
O_{A,B}(F(X))+O_{A,B}(1)=o\left(\frac X{\log X}\right).
\]
Thus
\[
\binom{\mathcal N_A(X)}2
\le
\frac12\left(\int_0^X f(t)\,dt\right)^2
+
o\left(\frac X Y\right).
\tag{14}\label{eq:small-tops-bound}
\]
By \eqref{eq:f-integral-asymptotic},
\[
\frac12\left(\int_0^X f(t)\,dt\right)^2
=
\frac{4\lambda X}{Y}
+
o\left(\frac X Y\right).
\tag{15}\label{eq:small-top-asymptotic}
\]

\proofheading{5. Large tops and the integer floor}

The only improvement over the continuous large-top majorant is that the integer count attached to
a top $t$ is bounded by $\lfloor G_X(t)\rfloor$, not merely by $G_X(t)$.
For $t\ge X$, define
\[
G_X(t):=\int_{t-X}^t f(s)\,ds.
\]
Since $f$ is nonincreasing, $G_X$ is also nonincreasing. Indeed, for $h>0$,
\[
G_X(t+h)
=
\int_{t-X}^t f(s+h)\,ds
\le
\int_{t-X}^t f(s)\,ds
=
G_X(t).
\]
Hence $\lfloor G_X(t)\rfloor$ is also nonincreasing. Since $G_X$ is monotone, it is measurable,
and therefore $\lfloor G_X\rfloor$ is measurable. On each compact interval $[X,M]$,
\[
0\le G_X(t)\le \int_0^X f(s)\,ds,
\]
because $t-X\ge0$ and $f$ is nonincreasing. Thus $G_X(t)$ is finite for every $t\in[X,M]$,
$\lfloor G_X\rfloor$ is bounded on $[X,M]$, and $f(t)\lfloor G_X(t)\rfloor$ is measurable and
integrable there.

Fix $t\in A$ with
\[
t>X+Z_0.
\]
Let
\[
r_t:=\#\{b\in A:t-X\le b<t\}.
\]
If $r_t=0$, then $r_t\le\lfloor G_X(t)\rfloor$ is trivial because $G_X(t)\ge0$. Assume
$r_t\ge1$, and list the contributing bottoms as
\[
b_1<b_2<\cdots<b_{r_t}<t,
\]
with
\[
b_i\in A,
\qquad
 t-X\le b_i<t.
\]
Because $t>X+Z_0$, each $b_i$ satisfies
\[
b_i\ge t-X>Z_0.
\]
Set $b_{r_t+1}:=t$. Then
\[
b_1<b_2<\cdots<b_{r_t}<b_{r_t+1}=t
\]
has exactly $r_t$ gaps, including the final gap ending at $t$. For every $1\le i\le r_t$, the
upper endpoint $b_{i+1}$ exceeds $Z_0$, so
\eqref{eq:gap-integral} applies to the gap $[b_i,b_{i+1}]$. Hence
\[
r_t
\le
\sum_{i=1}^{r_t}\int_{b_i}^{b_{i+1}} f(s)\,ds
=
\int_{b_1}^t f(s)\,ds
\le
\int_{t-X}^t f(s)\,ds
=
G_X(t).
\]
The quantity $G_X(t)$ is nonnegative, and $r_t$ is an integer. Hence the inequality $r_t\le
G_X(t)$ implies
\[
r_t\le \lfloor G_X(t)\rfloor.
\tag{16}\label{eq:floor-top-bound}
\]

Let
\[
t_1<t_2<\cdots<t_L
\]
be the elements of $A\cap(X+Z_0,M]$. If $L=0$, there is no contribution from this range. Assume
$L\ge1$. For $j\ge2$, the elements $t_{j-1}$ and $t_j$ are elements of $A$ with $t_{j-1}<t_j$
and
\[
t_j>X+Z_0>Z_0.
\]
The spacing estimate \eqref{eq:gap-integral}, which only requires two elements $y<z$ of $A$
with $z>Z_0$, therefore applies to the gap $[t_{j-1},t_j]$.

The next step converts the discrete sum over large tops into a gap-integral sum, losing only
the first large top. Because $\lfloor G_X(t)\rfloor$ is nonincreasing,
\[
\int_{t_{j-1}}^{t_j}
 f(s)\lfloor G_X(s)\rfloor\,ds
\ge
\lfloor G_X(t_j)\rfloor
\int_{t_{j-1}}^{t_j}f(s)\,ds
\ge
\lfloor G_X(t_j)\rfloor.
\]
Therefore, if $L\ge2$,
\[
\sum_{j=2}^L\lfloor G_X(t_j)\rfloor
\le
\int_{t_1}^{t_L} f(t)\lfloor G_X(t)\rfloor\,dt
\le
\int_X^M f(t)\lfloor G_X(t)\rfloor\,dt.
\]
Here extending the integral to all of $[X,M]$ is legitimate because
\[
f(t)\lfloor G_X(t)\rfloor\ge0.
\]
The first top $t_1$ is the only possible loss in this integral domination. Adding this
first-top loss gives
\[
\sum_{j=1}^L\lfloor G_X(t_j)\rfloor
\le
\lfloor G_X(t_1)\rfloor
+
\int_X^M f(t)\lfloor G_X(t)\rfloor\,dt.
\tag{17}\label{eq:integral-domination}
\]
For $L=1$, the same inequality is immediate because the integral term is nonnegative. Thus
\eqref{eq:integral-domination} holds for all $L\ge1$.

The endpoint term is negligible. Since $t_1>X$,
\[
\lfloor G_X(t_1)\rfloor
\le
G_X(t_1)
=
\int_{t_1-X}^{t_1}f(s)\,ds.
\]
Here $t_1-X>0$. Among intervals of length $X$ in $[0,\infty)$, the integral of a nonincreasing
function is maximized by $[0,X]$. Indeed, for $a\ge0$,
\[
\int_a^{a+X}f(s)\,ds
=
\int_0^X f(u+a)\,du
\le
\int_0^X f(u)\,du.
\]
Therefore, by \eqref{eq:f-integral-asymptotic},
\[
G_X(t_1)
\le
\int_0^X f(s)\,ds
=
O_{A,B}\left(\sqrt{\frac X Y}\right)
=
o\left(\frac X Y\right).
\tag{18}\label{eq:endpoint-loss}
\]
Consequently, if $L=0$, the estimate below is immediate because the left side is zero and the
integral is nonnegative. If $L\ge1$, \eqref{eq:floor-top-bound},
\eqref{eq:integral-domination}, and \eqref{eq:endpoint-loss} give the same estimate. Hence, in all
cases,
\[
\sum_{t\in A\cap(X+Z_0,M]} r_t
\le
\int_X^M f(t)\lfloor G_X(t)\rfloor\,dt
+
o\left(\frac X Y\right).
\]

The integer tops in the short interval $X<t\le X+Z_0$ also contribute only $o(X/Y)$. Indeed,
since $A\subseteq\N$ and $X$ and $Z_0$ are integers, the interval $(X,X+Z_0]\cap\N$ contains at
most $Z_0$ possible integer tops. This is the only use of the integrality of $Z_0$. Each such top has at most
$\mathcal N_A(X+Z_0)$ possible bottoms. By \eqref{eq:N-XplusZ0-bound}, the total contribution from these tops is
at most
\[
Z_0\mathcal N_A(X+Z_0)
=O_{A,B}\left(\sqrt{\frac X Y}\right)
=o\left(\frac X Y\right).
\]
Combining \eqref{eq:top-counting-identity}, \eqref{eq:small-tops-bound}, and the preceding
large-top bounds, we get
\[
X
\le
\frac12\left(\int_0^X f(t)\,dt\right)^2
+
\int_X^M f(t)\lfloor G_X(t)\rfloor\,dt
+
o\left(\frac X Y\right).
\tag{19}\label{eq:fundamental-upper-bound}
\]
It remains to evaluate the right-hand side sharply.

\proofheading{6. The continuous majorant without the floor}

Throughout parts 6 through 8,
\[
Y=\log X,
\]
with $X\to\infty$ through positive integers. Also $M=\Phi(X)$; as observed above,
$M\asymp_B X^2/Y>X$ for all sufficiently large $X$. Later $t=\Phi(r)$, $r=\psi(t)$, and
$q=X/r$, while $U$ and $Q$ are fixed real truncation parameters before the limit
$X\to\infty$ is taken. The lower-endpoint expansion feeds both the square term and the
large-$q$ floor tail; local Lipschitz control feeds the correction-tail estimate; and the
floor-saving main term uses the fixed-$Q$ change of variables together with
$G_X(\Phi(X/q))=q+o(1)$ uniformly on fixed $q$-ranges. For tail estimates below that are later
used after a truncation parameter is sent to infinity, implicit constants are independent of that
truncation parameter, while lower $X$-thresholds may depend on its fixed value. Dependence on a
fixed truncation parameter is displayed in subscripts such as $o_{A,B,U}$ or $o_{A,B,Q}$.
Define
\[
J_B(X)
:=
\frac12\left(\int_0^X f(t)\,dt\right)^2
+
\int_X^M f(t)G_X(t)\,dt.
\]
We prove
\[
J_B(X)
=
X+
\frac{\lambda\bigl(3-2B-2\log(\log 2)\bigr)X}{Y}
+
o\left(\frac X Y\right).
\tag{20}\label{eq:continuous-majorant-asymptotic}
\]
The proof of this asymptotic is split into the square term, the correction term on a fixed
truncation range, and a correction-tail estimate whose normalized constant is independent of the
truncation parameter.

The small-top term is already known from \eqref{eq:small-top-asymptotic}:
\[
\frac12\left(\int_0^X f(t)\,dt\right)^2
=
\frac{4\lambda X}{Y}
+
o\left(\frac X Y\right).
\tag{21}\label{eq:small-top-asymptotic-restated}
\]
Next,
\[
\int_X^M f(t)G_X(t)\,dt
=
X\int_X^M f(t)^2\,dt
+
R(X),
\tag{22}\label{eq:square-correction-split}
\]
where
\[
R(X):=
\int_X^M
f(t)\int_{t-X}^t\bigl(f(s)-f(t)\bigr)\,ds\,dt.
\]
Because $f$ is nonincreasing and $s\le t$ inside the inner integral, $R(X)\ge0$.

\proofsubheading{6.1. The square term and constant computation}

\begin{lemma}[Lower endpoint expansion]
As $X\to\infty$ through positive integers, let
\[
r_0=r_0(X):=\psi(X),
\qquad
w_0=w_0(X):=\log r_0.
\]
Then
\[
w_0=
\frac Y2+\frac12\log Y-\frac12\log(2\lambda)+o(1).
\tag{23}\label{eq:lower-endpoint-asymptotic}
\]
Equivalently,
\[
r_0=\psi(X)=\sqrt{\frac{XY}{2\lambda}}\,(1+o(1)).
\]
Consequently, for every fixed $Q>1$,
\[
r_0=\psi(X)=o(X/Q).
\]
\end{lemma}

\begin{proof}
For all sufficiently large $X$, we have $X>T_*$, so $r_0=\psi(X)$ is evaluated on the fixed
inverse branch. The defining identity is
\[
\Phi(e^{w_0})=\Phi(r_0)=X=e^Y,
\]
or equivalently
\[
X=\lambda\frac{e^{2w_0}}{H(w_0)}.
\]
Thus
\[
2w_0=Y+\log H(w_0)-\log\lambda.
\]
Since $X\to\infty$ and $\Phi(e^w)\to\infty$ on the chosen tail, we have $w_0\to\infty$. Also,
$H(w)=w+O_B(\log w)$, so $\log H(w_0)=o(w_0)$. The identity above gives $Y=2w_0+o(w_0)$, and
hence
\[
w_0\sim\frac Y2.
\]
Consequently,
\[
\log H(w_0)=O_B(\log Y),
\]
and the preceding identity improves to
\[
w_0=\frac Y2+O_B(\log Y).
\]
Therefore
\[
H(w_0)=\frac Y2+O_B(\log Y),
\]
and so
\[
\log H(w_0)=\log Y-\log 2+o(1).
\]
The fixed constant $B$ enters only inside the $O_B(\log Y)$ term before this logarithm; after the
logarithm it contributes no constant term. Substituting the last estimate into
\[
2w_0=Y+\log H(w_0)-\log\lambda
\]
gives \eqref{eq:lower-endpoint-asymptotic}. Exponentiating also gives the explicit endpoint
asymptotic
\[
r_0=\psi(X)=\sqrt{\frac{XY}{2\lambda}}\,(1+o(1)).
\]
For fixed $Q$, the stated consequence follows because
\[
\frac{r_0}{X/Q}
=
Q\sqrt{\frac{Y}{2\lambda X}}\,(1+o(1))\to0.
\]
\end{proof}

\begin{lemma}[Square-term asymptotic and constant computation]
As $X\to\infty$ through positive integers,
\[
X\int_X^M f(t)^2\,dt
=
X+
\frac{\lambda\bigl(1-2B+2\log(\lambda/2)\bigr)X}{Y}
+
o\left(\frac X Y\right).
\]
\end{lemma}

\begin{proof}
This is where the $1/Y$ constant in the continuous majorant is determined.

As observed in Section 3, $M=\Phi(X)\asymp_B X^2/Y$, so $M>X$ for all sufficiently large $X$.
For all sufficiently large $X$, we also have $X>T_*$ and $X\ge N_0$. Hence $\psi(X)\ge N_0$, and
throughout the range $X\le t\le M$, every argument of $\psi$ lies on the inverse
$\psi:[T_*,\infty)\to[N_0,\infty)$ of
$\Phi:[N_0,\infty)\to[T_*,\infty)$. This is the only branch used in all changes of variables
below.

Use the change of variables
\[
t=\Phi(e^w)=\lambda\frac{e^{2w}}{H(w)}.
\]
Then
\[
\psi(t)=e^w,
\qquad
f(t)=e^{-w}.
\]
Also,
\[
dt=\lambda e^{2w}\left(\frac2{H(w)}-\frac{H'(w)}{H(w)^2}\right)\,dw.
\]
Therefore
\[
f(t)^2\,dt
=
\lambda\left(\frac2{H(w)}-\frac{H'(w)}{H(w)^2}\right)\,dw.
\tag{24}\label{eq:square-change-differential}
\]
The lower endpoint $t=X$ corresponds to $w=w_0=\log\psi(X)$. The upper endpoint satisfies
$M=\Phi(X)=\Phi(e^Y)$, so under $t=\Phi(e^w)$ it corresponds to $e^w=X$, hence to $w=Y$. By
\eqref{eq:lower-endpoint-asymptotic}, $w_0<Y$ for all sufficiently large $X$, so the interval
$[w_0,Y]$ is ordered as written. Since
\[
\left(\frac1{H(w)}\right)'=-\frac{H'(w)}{H(w)^2},
\]
\eqref{eq:square-change-differential} gives
\[
\int_X^M f(t)^2\,dt
=
\lambda\left[
2\int_{w_0}^Y\frac{dw}{H(w)}
+
\frac1{H(Y)}-
\frac1{H(w_0)}
\right].
\tag{25}\label{eq:square-integral-exact}
\]
The estimate \eqref{eq:lower-endpoint-asymptotic} gives $w_0\sim Y/2$. Hence the interval
$[w_0,Y]$ has length $O_B(Y)$, and uniformly for $w\in[w_0,Y]$, we have $w\asymp Y$. Also,
\[
H(w)=w\left(1-\frac{\log w-B}{w}\right).
\]
Since
\[
\frac{\log w-B}{w}=O_B\left(\frac{\log Y}{Y}\right)
\]
uniformly on $[w_0,Y]$, after increasing the lower threshold in $X$ we may assume
\[
\left|\frac{\log w-B}{w}\right|\le \frac12
\qquad(w\in[w_0,Y]).
\]
The geometric expansion gives
\[
\frac1{H(w)}
=
\frac1w+
\frac{\log w-B}{w^2}
+
O_B\left(\frac{(\log Y)^2}{Y^3}\right).
\tag{26}\label{eq:H-inverse-expansion}
\]
The error term in \eqref{eq:H-inverse-expansion}, integrated over an interval of length $O_B(Y)$,
contributes
\[
O_B\left(\frac{(\log Y)^2}{Y^2}\right)=o\left(\frac1Y\right).
\]
Write
\[
w_0=\frac Y2+r_X.
\]
By \eqref{eq:lower-endpoint-asymptotic},
\[
r_X=
\frac12\log Y-\frac12\log(2\lambda)+o(1).
\tag{27}\label{eq:rX-asymptotic}
\]
This moving lower endpoint is the source of the $-2\log Y$ contribution that cancels the later
$+2\log Y$ term in the square-term constant computation.
Since $r_X=O_B(\log Y)$,
\[
w_0=\frac Y2\left(1+\frac{2r_X}{Y}\right),
\qquad
\frac{r_X^2}{Y^2}=O\left(\frac{(\log Y)^2}{Y^2}\right)=o\left(\frac1Y\right).
\]
Therefore Taylor expansion of $\log(1+2r_X/Y)$ gives
\[
\log w_0=
\log\frac Y2+
\frac{2r_X}{Y}+
o\left(\frac1Y\right),
\]
and hence
\[
2\int_{w_0}^Y\frac{dw}{w}
=
2\log 2-
\frac{4r_X}{Y}+
o\left(\frac1Y\right).
\]
Also,
\[
\int \frac{\log w-B}{w^2}\,dw
=
-\frac{\log w+1-B}{w}.
\]
Thus
\[
2\int_{w_0}^Y\frac{\log w-B}{w^2}\,dw
=
2\left[
-\frac{\log Y+1-B}{Y}
+
\frac{\log w_0+1-B}{w_0}
\right].
\]
Using $w_0=Y/2+r_X$ and $r_X=O_B(\log Y)$, the errors caused by $r_X$ in the denominator and
logarithm are
\[
O\left(\frac{(\log Y)^2}{Y^2}\right)=o\left(\frac1Y\right).
\]
Therefore
\[
\frac{\log w_0+1-B}{w_0}
=
\frac{2(\log Y-\log 2+1-B)}{Y}
+
o\left(\frac1Y\right).
\]
Hence
\[
2\int_{w_0}^Y\frac{\log w-B}{w^2}\,dw
=
\frac{2\log Y-4\log 2+2-2B}{Y}
+
o\left(\frac1Y\right).
\]
Finally,
\[
H(Y)=Y+O_B(\log Y),
\qquad
H(w_0)=\frac Y2+O(\log Y).
\]
Thus
\[
\frac1{H(Y)}=\frac1Y+O\left(\frac{\log Y}{Y^2}\right),
\qquad
\frac1{H(w_0)}=\frac2Y+O\left(\frac{\log Y}{Y^2}\right).
\]
Hence
\[
\frac1{H(Y)}-
\frac1{H(w_0)}
=
-\frac1Y+
o\left(\frac1Y\right).
\tag{28}\label{eq:endpoint-denominator-difference}
\]
\noindent\textit{Square-term estimate.} Combining \eqref{eq:square-integral-exact} through
\eqref{eq:endpoint-denominator-difference}, the three numerators of the $1/Y$ contributions inside
the square brackets, before substituting
the asymptotic for $r_X$, are respectively
\[
\begin{aligned}
2\int_{w_0}^Y \frac{dw}{w} &: -4r_X,\\
2\int_{w_0}^Y \frac{\log w-B}{w^2}\,dw
&: 2\log Y-4\log 2+2-2B,\\
\frac1{H(Y)}-\frac1{H(w_0)} &: -1.
\end{aligned}
\]
Equivalently, \eqref{eq:rX-asymptotic} gives
\[
-4r_X=-2\log Y+2\log(2\lambda)+o(1),
\]
which is the term cancelling the positive $2\log Y$ contribution below. Therefore, after using
\eqref{eq:rX-asymptotic}, the resulting coefficient is
\[
\begin{aligned}
&\bigl(-4r_X\bigr)
 +\bigl(2\log Y-4\log 2+2-2B\bigr)-1 \\
&=\bigl(-2\log Y+2\log(2\lambda)\bigr)
  +\bigl(2\log Y-4\log 2+2-2B\bigr)-1+o(1) \\
&=1-2B+2\log(2\lambda)-4\log 2+o(1) \\
&=1-2B+2\log(\lambda/2)+o(1).
\end{aligned}
\]
Thus the logarithmic terms cancel explicitly; the last line uses
$2\log(2\lambda)-4\log 2=2\log(\lambda/2)$. Also
$2\lambda\log 2=1$ and $\lambda/2=(4\log 2)^{-1}$. We obtain
\[
\int_X^M f(t)^2\,dt
=
\lambda\left[
2\log 2+
\frac{1-2B+2\log(\lambda/2)}{Y}
\right]
+
o\left(\frac1Y\right).
\]
Since $2\lambda\log 2=1$, this becomes
\[
\int_X^M f(t)^2\,dt
=
1+
\frac{\lambda\bigl(1-2B+2\log(\lambda/2)\bigr)}{Y}
+
o\left(\frac1Y\right).
\]
Thus
\[
\begin{aligned}
X\int_X^M f(t)^2\,dt
&=
X+
\frac{\lambda\bigl(1-2B+2\log(\lambda/2)\bigr)X}{Y}
+
o\left(\frac X Y\right).
\end{aligned}
\tag{29}\label{eq:square-term-asymptotic}
\]
This is the square-term asymptotic used in the final coefficient comparison.

\noindent\textit{Square-term constant computation.} In the bracket in
\eqref{eq:square-integral-exact}, the only $1/Y$-scale contributions are
\[
\begin{aligned}
\text{lower-endpoint displacement} &: \quad -4r_X,\\
\text{integral of }(\log w-B)/w^2 &: \quad 2\log Y-4\log 2+2-2B,\\
\text{endpoint denominator difference} &: \quad -1.
\end{aligned}
\]
Using $r_X=\frac12\log Y-\frac12\log(2\lambda)+o(1)$, these add to
\[
1-2B+2\log(\lambda/2)+o(1),
\]
which is the coefficient in \eqref{eq:square-term-asymptotic}. No $1/Y$-scale term is
left in the error: the integrated remainder in \eqref{eq:H-inverse-expansion}, the
$r_X$-dependent denominator and logarithm errors, and the endpoint expansion errors are all
$o(1/Y)$ before multiplication by $X$.
\end{proof}

\proofsubheading{6.2. The correction term}

\proofsubheading{\texorpdfstring{$L^1$}{L1} rescaling lemma}

\begin{lemma}[\texorpdfstring{$L^1$}{L1} rescaling]
Define
\[
g_X(v):=\sqrt{\frac{XY}{2\lambda}}\,f(Xv)
\qquad(v>0).
\]
The value at $v=0$ is irrelevant to all integrals below. For every fixed real $U>0$,
\[
g_X\to \bigl(v\mapsto v^{-1/2}\bigr)
\qquad\text{in }L^1((0,U)).
\tag{30}\label{eq:gX-L1-convergence}
\]
\end{lemma}

\begin{proof}
Fix $0<\varepsilon<\min\{U,1\}$. The parameter $\varepsilon$ is fixed before the limit
$X\to\infty$ is taken; only after the limsup estimate is obtained do we let
$\varepsilon\downarrow0$. The convergence is uniform on every interval $[\varepsilon,U]$,
by \eqref{eq:f-asymptotic}. On $(0,\varepsilon)$,
\[
\int_0^\varepsilon g_X(v)\,dv
=
\sqrt{\frac{Y}{2\lambda X}}
\int_0^{\varepsilon X} f(s)\,ds.
\]
Choose a fixed $T_1\ge T_*$ large enough for \eqref{eq:f-upper-bound} and
\eqref{eq:elementary-integral-bound} to hold on $[T_1,\infty)$. The compact initial part
contributes
\[
\sqrt{\frac{Y}{2\lambda X}}
\int_0^{T_1} f(s)\,ds
=
O_{A,B}\left(\sqrt{\frac YX}\right)=o(1).
\]
For each fixed $\varepsilon>0$ and all sufficiently large $X$, say with $\varepsilon X\ge T_1$,
the remaining part satisfies, by \eqref{eq:f-upper-bound} and
\eqref{eq:elementary-integral-bound},
\[
\int_{T_1}^{\varepsilon X} f(s)\,ds
\ll_B
\int_{T_1}^{\varepsilon X}\frac{ds}{\sqrt{s\log s}}
\ll_B
\sqrt{\frac{\varepsilon X}{\log(\varepsilon X)}}.
\]
Consequently,
\[
\int_0^{\varepsilon X} f(s)\,ds
\ll_{A,B}
1+
\sqrt{\frac{\varepsilon X}{\log(\varepsilon X)}}.
\]
For fixed $\varepsilon>0$ and all sufficiently large $X$ depending on this fixed $\varepsilon$,
\[
\log(\varepsilon X)=Y+\log\varepsilon,
\qquad
\limsup_{X\to\infty}\sqrt{\frac{Y}{\log(\varepsilon X)}}=1.
\]
Thus, for fixed $0<\varepsilon\le \min\{U,1\}$, taking $X\to\infty$ first with a post-limsup
constant independent of both $\varepsilon$ and $U$ gives the small-$v$ uniform-integrability
estimate
\[
\limsup_{X\to\infty}
\int_0^\varepsilon g_X(v)\,dv
\le
C_{A,B}\sqrt\varepsilon.
\]
Together with
\[
\int_0^\varepsilon v^{-1/2}\,dv=2\sqrt\varepsilon,
\]
this yields
\[
\limsup_{X\to\infty}
\int_0^\varepsilon |g_X(v)-v^{-1/2}|\,dv
\ll_{A,B}
\sqrt\varepsilon,
\]
where the implicit constant is independent of $0<\varepsilon\le \min\{U,1\}$. The threshold in
$X$ may depend on the fixed value of $\varepsilon$; only afterward do we let
$\varepsilon\downarrow0$.

Combining this with uniform convergence on $[\varepsilon,U]$, we have
\[
\limsup_{X\to\infty}
\|g_X-v^{-1/2}\|_{L^1(0,U)}
\le
C_{A,B}\sqrt\varepsilon
+
\limsup_{X\to\infty}
\|g_X-v^{-1/2}\|_{L^1(\varepsilon,U)}.
\]
The second term is $0$, and then letting $\varepsilon\downarrow0$ proves
\eqref{eq:gX-L1-convergence}.
\end{proof}

\proofsubheading{Truncated correction term}

In this subsection the order of limits is fixed: first fix $U\ge2$, then let $X\to\infty$,
and only after the normalized estimates are obtained let $U\to\infty$. The parameter $U$ never
depends on $X$.

Fix a real number $U\ge2$, and define
\[
R_U(X)
:=
\int_X^{UX}
 f(t)\int_{t-X}^t\bigl(f(s)-f(t)\bigr)\,ds\,dt.
\]
Here $U$ is fixed as $X\to\infty$. For fixed $U$, one has $UX<M$ for all sufficiently large
$X$, because
\[
\frac{M}{UX}\asymp_{B,U}\frac{X}{Y}\to\infty.
\]
Let
\[
h(u):=2\left(1-\sqrt{1-\frac1u}\right)-\frac1u
\qquad(u\ge1).
\]
With
\[
t=X\alpha,
\qquad
s=X\beta,
\]
we get
\[
\frac{Y}{2\lambda X}R_U(X)
=
\int_1^U
 g_X(\alpha)\int_{\alpha-1}^{\alpha}
 \bigl(g_X(\beta)-g_X(\alpha)\bigr)\,d\beta\,d\alpha.
\]
On $[1,U]$, $g_X(\alpha)\to \alpha^{-1/2}$ uniformly, again by \eqref{eq:f-asymptotic}, and the
functions $g_X$ are uniformly bounded there for all sufficiently large $X$. Also,
\[
\beta^{-1/2}\in L^1(0,U).
\]
Define
\[
K(\varphi)(\alpha):=\int_{\alpha-1}^{\alpha} \varphi(\beta)\,d\beta.
\]
Endpoint values are irrelevant for these Lebesgue integrals; in particular the integrable
singularity of $\beta^{-1/2}$ at $\beta=0$ causes no ambiguity at $\alpha=1$.
Then
\[
\|K(\varphi)\|_{L^\infty(1,U)}
\le
\|\varphi\|_{L^1(0,U)}.
\]
By \eqref{eq:gX-L1-convergence},
\[
K(g_X)\to K(\beta^{-1/2})
\]
uniformly on $[1,U]$. Moreover, $K(\beta^{-1/2})(\alpha)$ is finite and continuous on $[1,U]$.
The $L^1(0,U)$ convergence also gives $\|g_X\|_{L^1(0,U)}=O_{A,B,U}(1)$, so $K(g_X)$ is uniformly
bounded on $[1,U]$; and $g_X$ is uniformly bounded on $[1,U]$ for all sufficiently large $X$.
Writing $g(\alpha)=\alpha^{-1/2}$, for $\alpha\in[1,U]$ we have
\[
\begin{aligned}
&|g_X(\alpha)K(g_X)(\alpha)-g(\alpha)K(g)(\alpha)|
\\
&\le
\|g_X-g\|_{L^\infty(1,U)}\|K(g_X)\|_{L^\infty(1,U)}
\\
&\quad+
\|g\|_{L^\infty(1,U)}\|K(g_X)-K(g)\|_{L^\infty(1,U)},
\end{aligned}
\]
and
\[
|g_X(\alpha)^2-g(\alpha)^2|
\le
(\|g_X\|_{L^\infty(1,U)}+\|g\|_{L^\infty(1,U)})
\|g_X-g\|_{L^\infty(1,U)}.
\]
Both right-hand sides tend to $0$. Hence
\[
g_X(\alpha)K(g_X)(\alpha)-g_X(\alpha)^2
\to
\alpha^{-1/2}K(\beta^{-1/2})(\alpha)-\alpha^{-1}
\]
uniformly on $[1,U]$, and therefore in $L^1(1,U)$. Consequently,
\[
\frac{Y}{2\lambda X}R_U(X)
\to
\int_1^U
\left[
 \alpha^{-1/2}\int_{\alpha-1}^{\alpha} \beta^{-1/2}\,d\beta
-
\frac1\alpha
\right]d\alpha.
\]
Since
\[
\alpha^{-1/2}\int_{\alpha-1}^{\alpha} \beta^{-1/2}\,d\beta
=
2\left(1-\sqrt{1-\frac1\alpha}\right),
\]
we get
\[
R_U(X)
=
\frac{2\lambda X}{Y}
\int_1^U
\left[
2\left(1-\sqrt{1-\frac1\alpha}\right)-\frac1\alpha
\right]d\alpha
+
o_{A,B,U}\left(\frac X Y\right).
\tag{31}\label{eq:RU-limit}
\]
The limiting integrand is $h(\alpha)$, and
\[
h(\alpha)
=
\frac1{4\alpha^2}+O\left(\frac1{\alpha^3}\right)
\qquad(\alpha\to\infty).
\]
At $\alpha=1$ it equals $1$, and it is continuous on $[1,2]$. Hence the corresponding improper
integral converges absolutely.

\proofsubheading{Correction-tail estimate with truncation-independent constant}

\begin{lemma}[Correction-tail estimate]
For every fixed real $U\ge2$,
\[
0\le R(X)-R_U(X)\ll_B\frac{X}{UY},
\tag{32}\label{eq:R-tail-bound}
\]
where the implicit constant is independent of $U$, while the lower threshold for $X$ may depend
on the fixed value of $U$. Consequently,
\[
\lim_{U\to\infty}\limsup_{X\to\infty}
\frac{Y}{2\lambda X}\bigl(R(X)-R_U(X)\bigr)=0.
\]
\end{lemma}

\begin{proof}
Take a fixed real number $U\ge2$. For fixed $U$, all sufficiently large $X$ satisfy $UX<M$.
If $t\ge UX$ and $s\in[t-X,t]$, then
\[
s\in[t/2,t].
\]
For sufficiently large $X$, this interval lies in the smooth tail $(T_*,\infty)$. The lower
threshold for $X$ in this paragraph may depend on the fixed value of $U$, but the constant in
\eqref{eq:local-lipschitz} with $c=1/2$ depends only on $B$. Applying
\eqref{eq:local-lipschitz} with $c=1/2$, we get
\[
0\le f(s)-f(t)
\ll_B
\frac{f(t)}t(t-s).
\]
Therefore
\[
\int_{t-X}^t\bigl(f(s)-f(t)\bigr)\,ds
\ll_B
\frac{f(t)X^2}{t}.
\]
Hence
\[
0\le R(X)-R_U(X)
\ll_B
X^2\int_{UX}^M\frac{f(t)^2}{t}\,dt.
\]
By \eqref{eq:f-asymptotic},
\[
f(t)^2\ll_B\frac1{t\log t},
\]
and for $t\ge UX$ with $U\ge2$ we have $\log t\ge\log(UX)\ge Y$. Thus, enlarging the
upper limit to infinity, uniformly in $U\ge2$ in the constant, and with only the lower threshold
in $X$ allowed to depend on the fixed $U$,
\[
R(X)-R_U(X)
\ll_B
X^2\int_{UX}^{\infty}\frac{dt}{t^2\log t}.
\]
Moreover,
\[
\int_{UX}^{\infty}\frac{dt}{t^2\log t}
\le
\frac1Y\int_{UX}^{\infty}\frac{dt}{t^2}
=
\frac1{UXY}.
\]
Therefore
\[
R(X)-R_U(X)
\ll_B
\frac{X}{UY}.
\]
This proves \eqref{eq:R-tail-bound}. Equivalently, for every fixed $U\ge2$,
\[
\limsup_{X\to\infty}
\frac{Y}{2\lambda X}\bigl(R(X)-R_U(X)\bigr)
\le
\frac{C_B}{U},
\]
where the constant is independent of $U$. Hence
\[
\lim_{U\to\infty}\limsup_{X\to\infty}
\frac{Y}{2\lambda X}\bigl(R(X)-R_U(X)\bigr)=0.
\]
Here the threshold for $X$ may depend on the fixed value of $U$, while the implicit constant in
\eqref{eq:R-tail-bound} does not.
\end{proof}

From \eqref{eq:RU-limit} and the preceding tail estimate, for every fixed $U\ge2$,
\[
\limsup_{X\to\infty}
\left|
\frac{Y}{2\lambda X}R(X)
-
\int_1^\infty h(u)\,du
\right|
\le
\frac{C_B}{U}
+
\int_U^\infty |h(u)|\,du.
\]
Here we used $R(X)-R_U(X)\ge0$ and the normalized tail bound
\eqref{eq:R-tail-bound}; no sign information on $h$ is needed because
$\int_U^\infty |h(u)|\,du\to0$. Since $h(u)=O(u^{-2})$, the right-hand side tends to $0$
as $U\to\infty$. Therefore
\[
R(X)
=
\frac{2\lambda X}{Y}
\int_1^\infty
\left[
2\left(1-\sqrt{1-\frac1u}\right)-\frac1u
\right]du
+
o\left(\frac X Y\right).
\]
With
\[
z=\sqrt{1-\frac1u},
\]
one has
\[
u=\frac1{1-z^2},
\qquad
 du=\frac{2z}{(1-z^2)^2}\,dz.
\]
Moreover,
\[
h(u)\,du=(1-z)^2\frac{2z}{(1-z^2)^2}\,dz=\frac{2z}{(1+z)^2}\,dz.
\]
The integral becomes
\[
\int_0^1\frac{2z}{(1+z)^2}\,dz.
\]
Since
\[
\int\frac{2z}{(1+z)^2}\,dz
=
2\log(1+z)+\frac{2}{1+z},
\]
evaluation at $0$ and $1$ gives
\[
\int_0^1\frac{2z}{(1+z)^2}\,dz=2\log 2-1.
\]
Therefore
\[
R(X)
=
\frac{2\lambda(2\log 2-1)X}{Y}
+
o\left(\frac X Y\right).
\tag{33}\label{eq:R-asymptotic}
\]

\proofsubheading{6.3. Assembling \texorpdfstring{$J_B(X)$}{J-B(X)}}

The continuous-majorant coefficient now has three sources only: the small-top square
term \eqref{eq:small-top-asymptotic-restated}, the square term \eqref{eq:square-term-asymptotic},
and the correction term \eqref{eq:R-asymptotic}. Keeping these sources separate, from
\eqref{eq:small-top-asymptotic-restated}, \eqref{eq:square-correction-split},
\eqref{eq:square-term-asymptotic}, and \eqref{eq:R-asymptotic},
\[
\begin{aligned}
J_B(X)
&=
\frac{4\lambda X}{Y}
+
X
+
\frac{\lambda\bigl(1-2B+2\log(\lambda/2)\bigr)X}{Y} \\
&\qquad
+
\frac{2\lambda(2\log 2-1)X}{Y}
+
o\left(\frac X Y\right).
\end{aligned}
\]
Thus
\[
J_B(X)
=
X+
\frac{\lambda X}{Y}
\left[4+1-2B+2\log(\lambda/2)+2(2\log 2-1)\right]
+
o\left(\frac X Y\right).
\]
The bracket equals
\[
3-2B+4\log 2+2\log(\lambda/2).
\]
Since
\[
\lambda=\frac1{2\log 2},
\qquad
\frac{\lambda}{2}=\frac1{4\log 2},
\]
we have
\[
2\log(\lambda/2)=-4\log 2-2\log(\log 2),
\]
and therefore
\[
4\log 2+2\log(\lambda/2)=-2\log(\log 2).
\]
Therefore
\[
J_B(X)
=
X+
\frac{\lambda\bigl(3-2B-2\log(\log 2)\bigr)X}{Y}
+
o\left(\frac X Y\right),
\]
which proves \eqref{eq:continuous-majorant-asymptotic}.

\proofheading{7. The floor saving}

Define
\[
E(X):=\int_X^M f(t)\{G_X(t)\}\,dt,
\]
where
\[
\{x\}:=x-\lfloor x\rfloor.
\]
Since $0\le\{G_X(t)\}<1$, we have the trivial bounds
\[
0\le E(X)\le \int_X^M f(t)\,dt.
\]
Then
\[
\int_X^M f(t)\lfloor G_X(t)\rfloor\,dt
=
\int_X^M f(t)G_X(t)\,dt
-
E(X).
\tag{34}\label{eq:floor-decomposition}
\]

The floor saving is the following lemma.
\begin{lemma}[Floor saving]
With the definitions above,
\[
E(X)
=
2\lambda(1-\gamma)\frac X{\log X}
+o\left(\frac X{\log X}\right).
\tag{35}\label{eq:floor-saving}
\]
\end{lemma}

\begin{proof}
The proof has four separate steps: a fixed-$Q$ change of variables, a uniform approximation for
$G_X(\Phi(X/q))$, removal of the fractional-part discontinuities on the fixed interval
$1\le q\le Q$, and a large-$q$ tail bound with a normalized constant independent of $Q$. The
relevant discontinuities are those of $q\mapsto\{q\}$ at the integers; this is why the fixed-$Q$
convergence below first removes small neighborhoods of the finitely many integers in $[1,Q]$.

Put
\[
q(t):=Xf(t)=\frac X{\psi(t)}.
\]
All occurrences of $\psi$ in this part are on the inverse branch fixed above; after fixing
$Q$, the lower threshold in $X$ is enlarged whenever necessary so that all displayed arguments
of $\Phi$ and $\psi$ lie on this branch. Since $\psi$ is increasing, $q(t)$ is decreasing in $t$;
larger $q$ corresponds to smaller $t$ in the large-top interval.
Equivalently, if $r=\psi(t)$, then
\[
q=\frac Xr,
\qquad
 t=\Phi(r)=\Phi(X/q).
\]
Fix a real number $Q>1$. In this part, $Q$ is fixed before $X\to\infty$. All lower bounds on
$X$ may depend on $Q$, and all $o_{A,B,Q}$-terms may depend on $A,B,Q$ and the fixed threshold
data, but not on $X$. Only after the limiting sandwich is obtained do we let $Q\to\infty$.

For fixed $Q>1$, we have
\[
\Phi(X/Q)\asymp_{B,Q}\frac{X^2}{Y}>X
\]
for all sufficiently large $X$. Since $Q>1$ and $\Phi$ is increasing on the chosen branch,
\[
X<\Phi(X/Q)<M=\Phi(X)
\]
for all sufficiently large $X$. Thus $G_X(\Phi(X/q))$ is defined for every $q\in[1,Q]$ once
$X$ is large enough.

\proofsubheading{\texorpdfstring{Step 1: fixed-$Q$ $q$-change of variables}{Step 1: fixed-Q q-change of variables}}

\begin{lemma}[Fixed-$Q$ $q$-change of variables]
Here $Q>1$ is fixed before $X\to\infty$; no assertion is made with $Q=Q(X)$. For every bounded
measurable function $\Theta_X:[1,Q]\to\mathbb R$ satisfying $|\Theta_X|\le1$,
\[
\int_{\Phi(X/Q)}^M f(t)\Theta_X(q(t))\,dt
=
\frac{2\lambda X}{Y}
\int_1^Q\frac{\Theta_X(q)}{q^2}\,dq
+
o_{A,B,Q}\left(\frac X Y\right),
\tag{36}\label{eq:q-change}
\]
uniformly over all such measurable choices of $\Theta_X$.
\end{lemma}

\begin{proof}
The uniformity is over the bounded measurable test functions for this fixed $Q$; it is not a
statement with $Q=Q(X)$. The function $\Theta_X$ is allowed to depend on $X$, and no continuity or
regularity of $\Theta_X$ is used; the error is controlled below by a supremum before integration.
For fixed $Q>1$, assume $X$ is so large that $X/Q\ge N_0$. On the interval
\[
\Phi(X/Q)\le t\le M=\Phi(X),
\]
the monotonicity of $\Phi$ on $[N_0,\infty)$ and the inverse relation between $\Phi$ and
$\psi$ give
\[
X/Q\le\psi(t)\le X,
\]
and therefore
\[
1\le q(t)\le Q.
\]
The monotone $C^1$ change-of-variables theorem applies to these bounded measurable integrands.
Changing variables $r=\psi(t)$, and then $q=X/r$, gives the endpoint correspondence
\[
t=\Phi(X/Q)\Longleftrightarrow q=Q,
\qquad
 t=M=\Phi(X)\Longleftrightarrow q=1.
\]
Hence
\[
\begin{aligned}
\int_{\Phi(X/Q)}^M f(t)\Theta_X(q(t))\,dt
&=
\int_{X/Q}^X
\frac1r \Theta_X(X/r)\Phi'(r)\,dr \\
&=
\int_Q^1
\Theta_X(q)
\frac{\Phi'(X/q)}{X/q}
\left(-\frac{X}{q^2}\right)dq \\
&=
\int_1^Q
\Theta_X(q)
\frac{\Phi'(X/q)}{X/q}
\frac{X}{q^2}\,dq.
\end{aligned}
\]
We keep the product in the final integrand in this form because the first factor is uniformly
approximated by $2\lambda/Y+O_{B,Q}(\log Y/Y^2)$.
With $Q$ fixed, for $1\le q\le Q$, put
\[
w=\log(X/q)=Y-\log q.
\]
Then uniformly for $1\le q\le Q$,
\[
H(w)=w-\log w+B
=Y-\log Y+B-\log q+O_Q\left(\frac1Y\right)
=Y+O_{B,Q}(\log Y),
\]
and
\[
H'(w)=1-\frac1w=1+O_Q\left(\frac1Y\right).
\]
Therefore, uniformly for $1\le q\le Q$,
\[
\frac{\Phi'(X/q)}{X/q}
=
\lambda\frac{2H(w)-H'(w)}{H(w)^2}
=
\frac{2\lambda}{Y}
+
O_{B,Q}\left(\frac{\log Y}{Y^2}\right).
\]
In particular,
\[
\sup_{1\le q\le Q}
\left|
\frac{\Phi'(X/q)}{X/q}-\frac{2\lambda}{Y}
\right|
=
o_{A,B,Q}\left(\frac1Y\right).
\]
Thus the total error in replacing $\Phi'(X/q)/(X/q)$ by $2\lambda/Y$ is bounded by
\[
X
\sup_{1\le q\le Q}
\left|
\frac{\Phi'(X/q)}{X/q}-\frac{2\lambda}{Y}
\right|
\int_1^Q\frac{dq}{q^2}
=
o_{A,B,Q}\left(\frac XY\right),
\]
uniformly over all measurable $\Theta_X$ with $|\Theta_X|\le1$. The constant is uniform because $Q$ is
fixed before the limit $X\to\infty$; no uniformity in growing $Q$ is asserted here. This proves
\eqref{eq:q-change}.
\end{proof}

\begin{lemma}[Uniform approximation to $q$]
Uniformly for $1\le q\le Q$,
\[
G_X(\Phi(X/q))=q+O_{B,Q}\left(\frac YX\right).
\tag{37}\label{eq:GX-q-uniform}
\]
\end{lemma}

\begin{proof}
Put
\[
t=\Phi(X/q),
\qquad
1\le q\le Q.
\]
Then, since $H(Y-\log q)=Y+O_{B,Q}(\log Y)$,
\[
t=\lambda\frac{(X/q)^2}{H(Y-\log q)}
\asymp_{B,Q}\frac{X^2}{q^2Y},
\qquad
\frac Xt=O_B\left(\frac{q^2Y}{X}\right)=O_{B,Q}\left(\frac YX\right),
\]
and also
\[
\frac{qX}{t}=O_B\left(\frac{q^3Y}{X}\right)
=O_{B,Q}\left(\frac YX\right).
\]
In particular $t\to\infty$ uniformly for $1\le q\le Q$, so the threshold in
\eqref{eq:local-lipschitz} with $c=1/2$ is exceeded uniformly for fixed $Q$. For
$s\in[t-X,t]$, we have $s\in[t/2,t]$ for all sufficiently large $X$. By
\eqref{eq:local-lipschitz},
\[
f(s)=f(t)+O_B\left(\frac{f(t)}tX\right)
\]
uniformly in $q\in[1,Q]$. Therefore
\[
G_X(t)
=
\int_{t-X}^t f(s)\,ds
=
Xf(t)+O_B\left(\frac{f(t)X^2}{t}\right).
\]
Since
\[
\psi(\Phi(X/q))=X/q,
\qquad
 f(\Phi(X/q))=\frac qX,
\]
we obtain, using the preceding bound on $qX/t$,
\[
G_X(t)=q+O_B\left(\frac{qX}{t}\right)
=q+O_B\left(\frac{q^3Y}{X}\right)
=q+O_{B,Q}\left(\frac YX\right),
\]
which proves \eqref{eq:GX-q-uniform}.
\end{proof}

\proofsubheading{Step 2: fractional-part convergence away from discontinuities}

\begin{lemma}[Fractional-part convergence on a fixed $q$-range]
For every fixed real $Q>1$,
\[
\int_{\Phi(X/Q)}^M f(t)\{G_X(t)\}\,dt
=
\frac{2\lambda X}{Y}
\int_1^Q\frac{\{q\}}{q^2}\,dq
+
o_{A,B,Q}\left(\frac X Y\right).
\tag{38}\label{eq:floor-main-truncation}
\]
\end{lemma}

\begin{proof}
The following compactness argument is the discontinuity-removal step for the fractional-part map.
Since $q\mapsto\Phi(X/q)$ is continuous on $[1,Q]$ and $f$ is locally integrable, the primitive
\[
P(t):=\int_0^t f(s)\,ds
\]
is absolutely continuous on compact intervals, and hence continuous. For $t\ge X$,
\[
G_X(t)=P(t)-P(t-X),
\]
so $G_X$ is continuous on $[X,\infty)$. Hence the function
\[
\Theta_X(q):=\{G_X(\Phi(X/q))\}
\]
is measurable. Since \eqref{eq:q-change} is uniform over all measurable $|\Theta_X|\le1$, it may be
applied to this $X$-dependent choice.

The only obstruction to replacing $\{G_X(\Phi(X/q))\}$ by $\{q\}$ is crossing an integer; the next
paragraph removes small neighborhoods of the finitely many relevant integers.
For fixed $Q>1$ and $\eta>0$, choose $\mathcal U_\eta\subset[1,Q]$ to be a finite union of small
relatively open intervals around the integers in $[1,Q]$, using one-sided intervals at endpoint
integers if present, such that
\[
\int_{\mathcal U_\eta}\frac{dq}{q^2}<\eta.
\]
This is possible because there are only finitely many integers in $[1,Q]$. Endpoint behavior is
irrelevant for the Lebesgue integral; including endpoint integers only makes the uniform
convergence statement on the complement unambiguous. Let
\[
S_\eta:=[1,Q]\setminus \mathcal U_\eta.
\]
If $S_\eta=\varnothing$, the convergence claim below is vacuous. Otherwise, $S_\eta$ is compact
and has positive distance from the finite set $\mathbb Z\cap[0,Q+1]$; all other integers are
farther away from $[1,Q]$. Call this distance $\rho_\eta>0$. Thus
\[
\operatorname{dist}(q,\mathbb Z)\ge\rho_\eta
\qquad(q\in S_\eta).
\]
By the uniform convergence \eqref{eq:GX-q-uniform}, for large $X$,
\[
\left|G_X(\Phi(X/q))-q\right|<\frac{\rho_\eta}{2}
\]
uniformly for $q\in S_\eta$. Therefore no integer is crossed: $G_X(\Phi(X/q))$ and $q$ lie in
the same connected component of $\mathbb R\setminus\mathbb Z$. On the relevant component away
from the removed neighborhoods, the fractional-part map is Lipschitz with constant $1$, so
\[
\left|\{G_X(\Phi(X/q))\}-\{q\}\right|
\le
\left|G_X(\Phi(X/q))-q\right|.
\]
Hence
\[
\{G_X(\Phi(X/q))\}\to\{q\}
\]
uniformly for $q\in S_\eta$.

Applying \eqref{eq:q-change} to
\[
\Theta_X(q):=\{G_X(\Phi(X/q))\},
\]
and using that, for $t\in[\Phi(X/Q),M]$,
\[
\Phi(X/q(t))=\Phi(\psi(t))=t,
\]
we get
\[
\int_{\Phi(X/Q)}^M f(t)\{G_X(t)\}\,dt
=
\frac{2\lambda X}{Y}
\int_1^Q
\frac{\{G_X(\Phi(X/q))\}}{q^2}\,dq
+
o_{A,B,Q}\left(\frac X Y\right).
\]
On $S_\eta$, the integrand converges uniformly to $\{q\}/q^2$. On $\mathcal U_\eta$, both fractional
parts lie in $[0,1)$, so the total possible discrepancy is at most
\[
\int_{\mathcal U_\eta}\frac{dq}{q^2}<\eta.
\]
Therefore
\[
\limsup_{X\to\infty}
\left|
\int_1^Q
\frac{\{G_X(\Phi(X/q))\}-\{q\}}{q^2}\,dq
\right|
\le \eta.
\]
Letting $\eta\downarrow0$, we get \eqref{eq:floor-main-truncation}.
\end{proof}

\proofsubheading{\texorpdfstring{Step 3: large-$q$ floor-tail estimate with $Q$-independent constant}%
{Step 3: large-q floor-tail estimate with Q-independent constant}}

\begin{lemma}[Fixed-$Q$ large-$q$ floor-tail estimate]
For every fixed real $Q>1$, fixed independently of $X$, all sufficiently large $X$ satisfy
\[
\int_X^{\Phi(X/Q)} f(t)\{G_X(t)\}\,dt
\le
C_B\frac X{QY},
\tag{39}\label{eq:floor-tail-bound}
\]
where $C_B$ is independent of $Q$, while the lower threshold for $X$ may depend on the fixed
value of $Q$.
\end{lemma}

\begin{proof}
It remains to control the range
\[
X\le t\le\Phi(X/Q),
\]
which is the large-$q$ tail because $q(t)=X/\psi(t)$ decreases as $t$ increases. Splitting
$E(X)$ at the common endpoint
$t=\Phi(X/Q)$ changes no integral, since that endpoint has measure zero. Only a normalized
bound tending to zero as $Q\to\infty$ is
needed here; no asymptotic for this range is required. By the lower-endpoint expansion
\eqref{eq:lower-endpoint-asymptotic}, writing $r_0=\psi(X)$,
\[
\log r_0=\frac12Y+O_B(\log Y),
\]
and hence $X/r_0\to\infty$. Since $\psi$ is increasing, $q(t)=X/\psi(t)$ is decreasing, with
endpoint correspondence
\[
t=X\Longleftrightarrow q=\frac{X}{r_0},
\qquad
 t=\Phi(X/Q)\Longleftrightarrow q=Q.
\]
Thus this interval corresponds to $q\ge Q$ for all sufficiently large $X$. Endpoints do not
affect the integral. For each fixed $Q>1$, this interval is nonempty for all sufficiently large
$X$, because
\[
\Phi(X/Q)\asymp_{B,Q}\frac{X^2}{Y}\gg X.
\]
Since
\[
0\le\{G_X(t)\}<1,
\]
we have
\[
0\le
\int_X^{\Phi(X/Q)} f(t)\{G_X(t)\}\,dt
\le
\int_X^{\Phi(X/Q)} f(t)\,dt.
\]
For each fixed $Q>1$, the lower-endpoint lemma gives $r_0=\psi(X)=o(X/Q)$, so all sufficiently
large $X$ satisfy
\[
r_0<X/Q.
\]
This strict inequality fixes the orientation of the change of variables below. Here $X>T_*$ for
all sufficiently large $X$, so $r_0=\psi(X)$ is evaluated on the same inverse branch.

For fixed $Q$, all sufficiently large $X$ also satisfy $X/Q\ge N_0$, so both endpoints are on
the monotone branch of $\Phi$. Changing variables $t=\Phi(r)$, the last integral becomes
\[
\int_{r_0}^{X/Q}\frac{\Phi'(r)}r\,dr.
\]
Now, for all sufficiently large $r$,
\[
H(\log r)\asymp_B\log r,
\qquad
2H(\log r)-H'(\log r)\asymp_B\log r.
\]
Thus
\[
\frac{\Phi'(r)}r
=
\lambda\frac{2H(\log r)-H'(\log r)}{H(\log r)^2}
\ll_B
\frac1{\log r}
\]
uniformly for all $r\ge r_B$, with a constant depending only on $B$, after choosing a sufficiently
large threshold $r_B$. Increase the lower threshold in $X$ so that $r_0\ge r_B$ and
$X/Q\ge N_0$. Since
\[
\log r_0=\frac12Y+O_B(\log Y),
\]
we have $\log r_0\ge Y/3$ for all sufficiently large $X$. Throughout this range, therefore,
\[
\log r\ge\log r_0\ge \frac Y3.
\]
Consequently, using the interval orientation $r_0<X/Q$ just established,
\[
0\le \frac XQ-r_0\le \frac XQ.
\]
Therefore
\[
\int_{r_0}^{X/Q}\frac{\Phi'(r)}r\,dr
\ll_B
\int_{r_0}^{X/Q}\frac{dr}{Y}
\le
C_B\frac{X}{QY}.
\]
Therefore \eqref{eq:floor-tail-bound} follows. No sharper asymptotic is needed for this tail; the
$Q$-uniform normalized $O(1/Q)$ bound is enough after $Q\to\infty$. Equivalently,
\[
\frac{Y}{2\lambda X}
\int_X^{\Phi(X/Q)} f(t)\{G_X(t)\}\,dt
\ll_B
\frac1Q.
\tag{40}\label{eq:floor-tail-normalized}
\]
Using the normalized form \eqref{eq:floor-tail-normalized}, the estimate also gives, in
limsup form,
\[
\limsup_{X\to\infty}
\frac{Y}{2\lambda X}
\int_X^{\Phi(X/Q)} f(t)\{G_X(t)\}\,dt
\le
\frac{C_B}{Q},
\]
with $C_B$ independent of $Q$. Hence
\[
\lim_{Q\to\infty}\limsup_{X\to\infty}
\frac{Y}{2\lambda X}
\int_X^{\Phi(X/Q)} f(t)\{G_X(t)\}\,dt=0.
\]
\end{proof}

Combining \eqref{eq:floor-main-truncation} and \eqref{eq:floor-tail-bound}, and keeping the order
of limits $X\to\infty$ first with $Q$ fixed and then $Q\to\infty$, gives the desired sandwich.
The implicit constant in the normalized $O(1/Q)$ tail bound is independent of $Q$; only the lower
threshold for $X$ may depend on the fixed value of $Q$. This independence is the only property of
the large-$q$ tail used in the final sandwich. Let
\[
\mathcal I_Q:=\int_1^Q\frac{\{q\}}{q^2}\,dq,
\qquad
\mathcal I:=\int_1^\infty\frac{\{q\}}{q^2}\,dq.
\]
The integral defining $\mathcal I$ converges because
\[
0\le\frac{\{q\}}{q^2}\le\frac1{q^2}.
\]
For fixed $Q>1$, write
\[
E_{\le Q}(X):=\int_{\Phi(X/Q)}^M f(t)\{G_X(t)\}\,dt,
\qquad
E_{>Q}(X):=\int_X^{\Phi(X/Q)} f(t)\{G_X(t)\}\,dt.
\]
The subscripts refer to the corresponding $q$-range, not to the orientation of the
$t$-intervals.
Then $E(X)=E_{\le Q}(X)+E_{>Q}(X)$ up to the common endpoint, which has measure zero, and
$E_{>Q}(X)\ge0$. By \eqref{eq:floor-main-truncation},
\[
\frac{Y}{2\lambda X}E_{\le Q}(X)\to \mathcal I_Q,
\]
while the normalized tail estimate above, with a $Q$-independent constant, gives
\[
0\le \limsup_{X\to\infty}\frac{Y}{2\lambda X}E_{>Q}(X)\le\frac{C_B}{Q}.
\]
In unnormalized form, the compact summary is
\[
E_{\le Q}(X)=\frac{2\lambda X}{Y}\mathcal I_Q+o_{A,B,Q}\left(\frac X Y\right),
\qquad
0\le E_{>Q}(X)\ll_B \frac{X}{QY}
\]
with $X\to\infty$ at fixed $Q$, and only then $Q\to\infty$.
The order of limits in the next displays is
\[
X\to\infty\quad\text{with }Q\text{ fixed, and then}\quad Q\to\infty.
\]
For every fixed $Q>1$, the preceding estimates imply
\[
\liminf_{X\to\infty}
\frac{Y E(X)}{2\lambda X}
\ge \mathcal I_Q,
\qquad
\limsup_{X\to\infty}
\frac{Y E(X)}{2\lambda X}
\le \mathcal I_Q+\frac{C_B}{Q},
\]
with $C_B$ enlarged if necessary. Equivalently,
\[
\mathcal I_Q
\le
\liminf_{X\to\infty}
\frac{Y E(X)}{2\lambda X}
\le
\limsup_{X\to\infty}
\frac{Y E(X)}{2\lambda X}
\le
\mathcal I_Q+\frac{C_B}{Q}.
\]
Letting $Q\to\infty$, using $\mathcal I_Q\uparrow \mathcal I$ by monotone
convergence and $C_B/Q\to0$, gives the normalized convergence
\[
\frac{Y E(X)}{2\lambda X}\to \mathcal I.
\]
Equivalently,
\[
E(X)
=
\frac{2\lambda X}{Y}
\int_1^\infty\frac{\{q\}}{q^2}\,dq
+
o\left(\frac X Y\right).
\tag{41}\label{eq:E-asymptotic}
\]
It remains to compute the integral. Endpoint values of the fractional-part function are
immaterial here, since they affect only a measure-zero set. By nonnegativity, monotone
convergence permits the split into unit intervals:
\[
\int_1^\infty\frac{\{q\}}{q^2}\,dq
=
\sum_{m=1}^\infty
\int_m^{m+1}\frac{q-m}{q^2}\,dq.
\]
For each $m\ge1$,
\[
\int_m^{m+1}\frac{q-m}{q^2}\,dq
=
\log\frac{m+1}{m}-\frac1{m+1}.
\]
Therefore
\[
\sum_{m=1}^K
\left(
\log\frac{m+1}{m}-\frac1{m+1}
\right)
=
\log(K+1)-\sum_{m=1}^K\frac1{m+1}.
\]
Since
\[
\sum_{m=1}^K\frac1{m+1}
=
\sum_{j=1}^{K+1}\frac1j-1,
\]
we get
\[
\sum_{m=1}^K
\left(
\log\frac{m+1}{m}-\frac1{m+1}
\right)
=
\log(K+1)-\sum_{j=1}^{K+1}\frac1j+1.
\]
Letting $K\to\infty$, and using the definition of $\gamma$,
\[
\int_1^\infty\frac{\{q\}}{q^2}\,dq=1-\gamma.
\]
This value is positive because the integrand is nonnegative and positive on a set of positive
measure; numerically, $1-\gamma=0.42278\ldots$. Substituting this into
\eqref{eq:E-asymptotic} proves the floor-saving lemma, namely \eqref{eq:floor-saving}.
\end{proof}

\proofheading{8. Final contradiction}

We have proved the two asymptotic evaluations
\[
J_B(X)
=
X+
\frac{\lambda\bigl(3-2B-2\log(\log 2)\bigr)X}{Y}
+
o\left(\frac X Y\right),
\]
and
\[
E(X)=2\lambda(1-\gamma)\frac X Y+
o\left(\frac X Y\right).
\]
Recall that $X$ in this final comparison runs through sufficiently large positive integers.
The final comparison uses \eqref{eq:fundamental-upper-bound}, \eqref{eq:floor-decomposition},
\eqref{eq:continuous-majorant-asymptotic}, and \eqref{eq:floor-saving}; the continuous
asymptotic \eqref{eq:continuous-majorant-asymptotic} is assembled from
\eqref{eq:small-top-asymptotic-restated}, \eqref{eq:square-term-asymptotic}, and
\eqref{eq:R-asymptotic}. From \eqref{eq:fundamental-upper-bound},
\[
X
\le
\frac12\left(\int_0^X f(t)\,dt\right)^2
+
\int_X^M f(t)\lfloor G_X(t)\rfloor\,dt
+
o\left(\frac X Y\right).
\]
Using \eqref{eq:floor-decomposition} and the definition of $J_B(X)$, this becomes
\[
X\le J_B(X)-E(X)+o\left(\frac X Y\right).
\]
By \eqref{eq:continuous-majorant-asymptotic} and \eqref{eq:floor-saving},
\[
\begin{aligned}
X
&\le
X+
\frac{\lambda\bigl(3-2B-2\log(\log 2)\bigr)X}{Y} \\
&\qquad
-
\frac{2\lambda(1-\gamma)X}{Y}
+
o\left(\frac X Y\right).
\end{aligned}
\]
The coefficient of $X/Y$ is
\[
\lambda\bigl(3-2B-2\log(\log 2)\bigr)-2\lambda(1-\gamma)
=
-2\lambda(B-B_1).
\]
Let
\[
\delta:=B-B_1>0.
\]
Hence the preceding inequality gives
\[
0\le -\frac{2\lambda\delta X}{Y}+o\left(\frac X Y\right).
\]
After division by $X/Y$ along integer $X\to\infty$, write the error term as
$\eta_X=o(1)$. Thus
\[
0\le -2\lambda\delta+\eta_X.
\]
For all sufficiently large integer $X$, $|\eta_X|\le\lambda\delta$, and hence
\[
0\le -2\lambda\delta+\eta_X\le -\lambda\delta<0,
\]
which is impossible.

Thus the assumed eventual upper bound \eqref{eq:upper-bound-assumption} cannot hold. Given an
integer $N$, choose an integer $N'\ge \max\{N,2\}$. If every $n\ge N'$ satisfied the non-strict
upper bound, then the forbidden eventual upper bound would hold from $N'$ onward. Therefore there
exists $n\ge N'$ such that
\[
u(n)>\lambda\frac{n^2}{\log n-\log\log n+B}.
\]
This proves the main assertion.

\proofheading{Fixed-constant limsup consequence}

Fix a real constant $C$ and fix one value $B^*>B_1$. Apply the main assertion just proved with
this value of $B^*$. Then there is an unbounded sequence $n_j$ such that
\[
u(n_j)
>
\lambda\frac{n_j^2}{\log n_j-\log\log n_j+B^*}.
\]
This uses no limiting process $B^*\downarrow B_1$ inside the contradiction argument and does not
substitute $B=B_1$ in the pointwise inequality. Discarding finitely many terms if necessary, both
quantities
\[
\log n_j-\log\log n_j+C
\]
and
\[
\log n_j-\log\log n_j+B^*
\]
are positive along this sequence. Therefore the following division and multiplication preserve
the inequality sign, and along this unbounded sequence,
\[
\frac{u(n_j)(\log n_j-\log\log n_j+C)}{n_j^2}
>
\lambda
\frac{\log n_j-\log\log n_j+C}
{\log n_j-\log\log n_j+B^*}.
\]
The fixed additive difference $B^*-C$ is negligible compared with $\log n_j$; explicitly,
\[
\frac{\log n_j-\log\log n_j+C}
{\log n_j-\log\log n_j+B^*}
=
1-
\frac{B^*-C}{\log n_j-\log\log n_j+B^*}
\to1.
\]
The right-hand side tends to $\lambda$. Since the left-hand side is bounded below along this
sequence by quantities tending to $\lambda$, the full limsup is at least $\lambda$. Hence, for
every fixed real $C$,
\[
\limsup_{n\to\infty}
\frac{u(n)(\log n-\log\log n+C)}{n^2}
\ge\lambda.
\]
Since
\[
\lambda=\frac1{2\log 2},
\]
this is the fixed-constant limsup assertion in the theorem. Taking $C=B_1$ gives the endpoint
limsup version there. Since
\[
\frac{\log n}{\log n-\log\log n}\to1,
\]
the conclusion with $C=0$ also gives
\[
\limsup_{n\to\infty}
\frac{u(n)\log n}{n^2}
\ge
\frac1{2\log 2}.
\]
\end{proof}

\medskip
\phantomsection
\pdfbookmark[1]{Authorship, Lean proof, and verification disclosure}{disclosure}
\noindent\textit{Authorship, Lean proof, and verification disclosure.} GPT-5.5 Pro is listed as
first author because it produced nearly all of the mathematical derivation, proof structuring, and
LaTeX drafting through prompted interaction. Lech Mazur prompted, selected, audited, edited,
curated, and is responsible for the public posting of this draft. A companion Lean proof is
available at \url{https://github.com/lechmazur/erdos_1194/}. This public draft has not been
independently refereed, and neither the Lean proof nor the other checks should be regarded as
journal peer review.

\begin{thebibliography}{99}

\bibitem{AlexeevMixon2026}
B. Alexeev and D. G. Mixon,
\emph{Forbidden Sidon subsets of perfect difference sets, featuring a human-assisted proof},
arXiv:2510.19804v2 [math.CO], submitted 2025-10-22, revised 2026-01-16,
\url{https://doi.org/10.48550/arXiv.2510.19804}.

\bibitem{Bloom1194}
T. F. Bloom,
\emph{Erd\H{o}s Problem \#1194},
Erd\H{o}s Problems,
\url{https://www.erdosproblems.com/1194},
page last edited 2026-04-24, accessed 2026-04-29.

\bibitem{CillerueloNathanson2008}
J. Cilleruelo and M. B. Nathanson,
\emph{Perfect difference sets constructed from Sidon sets},
Combinatorica \textbf{28} (2008), 401--414,
\url{https://doi.org/10.1007/s00493-008-2339-4}.

\bibitem{Erdos1980}
P. Erd\H{o}s,
\emph{A survey of problems in combinatorial number theory},
Annals of Discrete Mathematics \textbf{6} (1980), 89--115,
\url{https://doi.org/10.1016/S0167-5060(08)70697-6}.

\bibitem{HalberstamRoth1966}
H. Halberstam and K. F. Roth,
\emph{Sequences. Vol. I},
Clarendon Press, Oxford, 1966.

\bibitem{Lev2004}
V. F. Lev,
\emph{Reconstructing integer sets from their representation functions},
Electronic Journal of Combinatorics \textbf{11} (2004), no.~1, Research Paper 78, 6 pp.,
\url{https://doi.org/10.37236/1831}.

\end{thebibliography}

\end{document}
